\section{Discriminant of a sesquilinear form}

Throughout this paragraph, A denotes a commutative ring, J an automorphism of A, and E a free A-module of finite dimension n.

DEFINITION 1. --- Given a form \(\Phi\) sesquilinear for J on E and a system \(S = (x_1, \ldots, x_n)\) of n elements of E, one calls the discriminant of \(\Phi\) with respect to this system, and denotes it by \(D_{\Phi}(x_1, \ldots, x_n)\) or \(D_{\Phi}(S)\) the element det \((\Phi(x_i, x_j))\) of A.

If \((e_{1},\ldots,e_{n})\) is a basis of E, the discriminant of \(\Phi\) with respect to this basis is none other than the determinant of the matrix of \(\Phi\) with respect to this basis.

It follows from the definition of the extension of \(\Phi\) to \(\wedge\) E ( \(\S\) 1, no. 9) that we have

\[
\mathrm{D} _ {\Phi} (x _ {1}, \dots , x _ {n}) = \Phi_ {(n)} (x _ {1} \wedge \dots \wedge x _ {n}, x _ {1} \wedge \dots \wedge x _ {n}),\tag{1}
\]

where \(\Phi_{(n)}\) denotes the extension of \(\Phi\) to \(\wedge\) E. For every permutation \(\sigma \in \mathfrak{S}_n\) , we therefore have

\[
\mathrm{D} _ {\Phi} (x _ {\sigma (\mathbf {1})}, \dots , x _ {\sigma (n)}) = \mathrm{D} _ {\Phi} (x _ {\mathbf {1}}, x _ {\mathbf {2}}, \dots , x _ {n}).
\]

Example. --- Let B be an algebra over the ring A, such that B is a free A-module of finite dimension n. Then the map \((x, y) \to \mathrm{Tr}_{\mathrm{B}/\mathrm{A}}(xy)\) (chap.~VIII, § 12, no. 2) is a bilinear form on B. Given a system \((x_{1}, \ldots, x_{n})\) of n elements of B, the discriminant of this form with respect to this system is called the discriminant of the system \((x_{1}, \ldots, x_{n})\) over A, and is denoted \(\mathrm{D}_{\mathrm{B}/\mathrm{A}}(x_{1}, \ldots, x_{n})\) . We thus have

\[
\mathrm{D} _ {\mathrm{B/A}} (x _ {1}, \dots , x _ {n}) = \det (\operatorname{Tr} _ {\mathrm{B/A}} (x _ {i} x _ {j})).\tag{2}
\]

Remark. --- Let \((e_{1},\ldots,e_{n})\) be a basis of B over A, and \(e_{i}e_{j}=\sum_{k=1}^{n}c_{ijk}e_{k}\;(c_{ijk}\in\mathrm{A})\) the multiplication table of this basis (chap.~II, § 7, no. 2). Since the matrix of the A-linear endomorphism \(x\to e_{k}x\) of B with respect to \((e_{r})\) is \((c_{ksr})\) , one has \(\mathrm{Tr}_{\mathrm{B}/\mathrm{A}}(e_{k})=\sum_{r=1}^{n}c_{krr}\) , whence \(\mathrm{Tr}_{\mathrm{B}/\mathrm{A}}(e_{i}e_{j})=\sum_{k,r}c_{ijk}c_{krr}\) . It follows that one has

\[
\mathrm{D} _ {\mathbf {B} / \mathbf {A}} (e _ {1}, \dots , e _ {n}) = \det _ {i, j} (\sum_ {k, r} c _ {i j k} c _ {k r r}).\tag{3}
\]

PROPOSITION 1. — Let \(\Phi\) be a sesquilinear form for J on E, \((x_1,\ldots,x_n)\) be a system of \(n\) elements of E, and \((a_{ij})_{(i,j=1,\ldots,n)}\) a family of \(n^2\) elements of A; set \(y_j = \sum_{i=1}^{n} a_{ji} x_i\) . We then have

\[
\mathrm{D} _ {\Phi} (y _ {1}, \dots , y _ {n}) = \det (a _ {i j}). \det (a _ {i j}) ^ {\mathrm{J}}. \mathrm{D} _ {\Phi} (x _ {1}, \dots , x _ {n}).
\]

Indeed, as \(\Phi\) is a sesquilinear form, we have \(\Phi(y_i, y_j) = \sum_{k,m} a_{ik} \Phi(x_k, x_m) a_{jm}^{\mathfrak{J}}\) . Therefore, if we denote by \(A\) the matrix \((a_{ij})\) , the matrix \((\Phi(y_i, y_j))\) is equal to \(A \cdot (\Phi(x_i, x_j)) \cdot {}^t A^{\mathfrak{J}}\) . Since \(\det(^t A) = \det(A)\) and \(\det(A^{\mathfrak{J}}) = \det(A)^{\mathfrak{J}}\) , our assertion is proved.

In particular, if \((e_{i})\) and \((e_{i}^{\prime})\) are two bases of E, D and D' the discriminants of \(\Phi\) with respect to these bases, and a the determinant of the change-of-basis matrix from the basis \((e_{i})\) to the basis \((e_{i}^{\prime})\) , one has

\[
\mathrm{D} ^ {\prime} = a a ^ {\mathrm{J}} \mathrm{D}.\tag{4}
\]

It follows from prop. 1 that, if \((e_i)\) is a basis of E and \((x_i)\) any arbitrary system of \(n\) elements of E, \(\mathrm{D}_{\Phi}(e_1,\ldots ,e_n)\) divides \(\mathrm{D}_{\Phi}(x_1,\dots,x_n)\) . In particular, the discriminants of \(\Phi\) with respect to any two bases of E generate the same principal ideal of A.

Let \((\mathbf{E}_i)_{i\in I}\) be a finite family of free A-modules of finite dimensions, \(\Phi_i\) a sesquilinear form for J on \(\mathbf{E}_i\) , and \(\mathbf{B}_i\) a basis of \(\mathbf{E}_i\) . If \(\Phi\) denotes the direct sum of \(\Phi_i\) ( \(\S 1\) , no. 3) and B the basis of \(\prod E_i\) obtained by taking the union of \(\mathbf{B}_i\) , we obviously have

\[
\mathrm{D} _ {\Phi} (\mathrm{B}) = \prod_ {i \in \mathrm{I}} \mathrm{D} _ {\Phi_ {i}} (\mathrm{B} _ {i}).\tag{5}
\]

Let \(\Phi\) be a sesquilinear form for \(J\) on \(E\) , \(h\) a homomorphism of \(A\) into a commutative ring \(A'\) , \(\Phi'\) the sesquilinear form on \(A' \otimes_A E\) obtained by extension of \(\Phi\) ( \(\S 1\) , no. 4) and \((x_1, \ldots, x_n)\) an arbitrary system of elements of \(E\) . Since \(A' \otimes_A E\) is an \(A'\) -free module, \(D_{\Phi'}(1 \otimes x_1, \ldots, 1 \otimes x_n)\) is defined, and we obviously have

\[
\mathrm{D} _ {\Phi^ {\prime}} (1 \otimes x _ {1}, \dots , 1 \otimes x _ {n}) = h (\mathrm{D} _ {\Phi} (x _ {1}, \dots , x _ {n})).\tag{6}
\]

Example. --- Let B be an algebra over A which is a free A-module of finite dimension \(n\) , \((x_{1},\ldots,x_{n})\) a basis of B over A, and m an ideal of A. If we denote by \(h\) the canonical homomorphism from B onto B/mB, \((h(x_{1}),\ldots,h(x_{n}))\) is a basis of B/mB over A/m (chap.~I, § 6, no. 5, prop. 5), and B/mB is isomorphic to \((A/m)\otimes_{A}B\) . We therefore have

\[
\mathrm{D} _ {(\mathrm{B} / \mathfrak {m B}) / (\mathrm{A} / \mathfrak {m})} (h (x _ {1}), \dots , h (x _ {n})) = h (\mathrm{D} _ {\mathrm{B} / \mathrm{A}} (x _ {1}, \dots , x _ {n})).
\]

PROPOSITION 2. --- Assume that A is an integral domain. Let \(\Phi\) be a sesquilinear form for J on E and \((e_1, \ldots, e_n)\) a basis of E, such that \(\mathrm{D}_{\Phi}(e_1, e_2, \ldots, e_n) \neq 0\) .

\begin{enumerate}
\def\labelenumi{\alph{enumi})}
\item
  For a system \((x_{1},\ldots ,x_{n})\) of \(n\) elements of E to be free, it is necessary and sufficient that \(\mathrm{D}_{\Phi}(x_1,\dots,x_n)\) be \(\neq 0\) .
\item
  For a system \((x_{1},\ldots ,x_{n})\) of \(n\) elements of E to be a basis of E, it is necessary and sufficient that \(\mathrm{D}_{\Phi}(x_1,\dots,x_n)\) and \(\mathrm{D}_{\Phi}(e_1,\dots,e_n)\) be associated elements in A (cf.~Chap.~VI, § 1, no.~5).
\end{enumerate}

Let us set \(x_j = \sum_{i=1}^{n} a_{ji} e_i\) ( \(a_{ji} \in A\) ). Let us first prove \(a\) ). If \(D_{\Phi}(x_1, \ldots, x_n) = 0\) , one has \(\det(a_{ji})\) . \(\det(a_{ji})^{\mathfrak{J}} = 0\) (prop. 1) since \(D_{\Phi}(e_1, \ldots, e_n) \neq 0\) and A is an integral domain; we therefore have \(\det(a_{ji}) = 0\) , and the vectors \(x_j\) are linearly dependent (chap.~III, § 7, no. 1, th. 1, applied to the vector space \(K \otimes_A E\) , where K denotes the field of fractions of A). Conversely, if these vectors are linearly dependent, we have \(\det(a_{ji}) = 0\) ( \(ibid.\) ), whence \(D_{\Phi}(x_1, \ldots, x_n) = 0\) (prop. 1).

We now prove \(b\) ). If \(\mathrm{D}_{\Phi}(x_1, \ldots, x_n)\) and \(\mathrm{D}_{\Phi}(e_1, \ldots, e_n)\) are associated in A, prop. 1 shows that \(\det(a_{ij}) \cdot \det(a_{ij})^{\mathrm{J}}\) is invertible in A. Thus \(\det(a_{ij})\) is also invertible in A; therefore the matrix \((a_{ij})\) on A is invertible (chap.~III, § 6, no. 5, th. 2), and the endomorphism \(g\) of E defined by \(g(e_i) = x_i\) ( \(i = 1, \ldots, n\) ) is an automorphism; consequently \((x_1, \ldots, x_n)\) is a basis of E. The converse follows immediately from Prop. 1.

PROPOSITION 3. --- Let \(\Phi\) be a sesquilinear form for J on E, and S a basis of E. The following conditions are equivalent:

\begin{enumerate}
\def\labelenumi{\alph{enumi})}
\item
  The map \(s_{\Phi}\) from E into \(E^{*}\) associated with \(\Phi\) is bijective.
\item
  The map \(d_{\Phi}\) from E into \(E^{*}\) associated with \(\Phi\) is bijective.
\item
  The element \(D_{\Phi}(S)\) is invertible in A.
\end{enumerate}

Indeed the condition \(c\) expresses that the matrix of \(\Phi\) with respect to S is invertible (chap. III, § 6, no. 5, th. 2). Therefore \(c\) is equivalent to \(a\) ( \(§ 1\) , no. 10); likewise \(c\) is equivalent to \(b\) ).

PROPOSITION 4. --- Assume A is an integral domain. Let S be a basis of E. A necessary and sufficient condition for a sesquilinear form \(\Phi\) on E to be nondegenerate is that one has D \(_{\Phi}\) (S) ≠ 0.

Indeed, let K be the field of fractions of A, and let \(\Phi'\) be the extension of \(\Phi\) to the K-vector space \(K\otimes_{\mathbf{A}}E\) ; we identify E with a subset of this vector space. The relation \(D_{\Phi}(S)\neq0\) is then equivalent to \(D_{\Phi'}(S)\neq0\) by formula (6), which itself expresses that \(s_{\Phi'}\) is bijective (prop. 3), that is, \(\Phi'\) is nondegenerate ( \(\S1\) , no. 6, prop. 6). Now, for every \(x\in K\otimes_{\mathbf{A}}E\) , there exists \(a\in A\) such that \(ax\in E\) . Therefore, for \(\Phi\) to be degenerate, it is necessary and sufficient that \(\Phi'\) be so. This proves our assertion.

PROPOSITION 5. --- Let A be a field, B a commutative algebra of finite dimension n over A, and S a basis of B. For B to be separable (chap. VIII, § 7, no. 5, def. 1) it is necessary and sufficient that one have D \(_{B/A}\) (S) ≠ 0.

Let A' be the algebraic closure of A, and let B' be the algebra A' ⊗\_A B over A'. If B is separable, B' is semisimple (Chap. VIII, § 7, no. 5, cor. of Prop. 7) and is therefore a direct product of n fields isomorphic to A' (Chap. VIII, § 6, no. 4, cor. of Prop. 9). If S' denotes the canonical basis of B' (identified with A'	extsuperscript{n}), we have D\_\{B'/A'\}(S') = 1, hence D\_\{B'/A'\}(S) ≠ 0 (prop. 1) and D\_\{B/A\}(S) ≠ 0 (formula (6)).

Conversely, suppose we have \(D_{B/A}(S) \neq 0\) . To show that B is separable, it suffices to show that \(B'\) is semisimple, that is, that it admits no nilpotent element \(\neq 0\) . Now, if \(x'\) were a nonzero nilpotent element of \(B'\) , one could take it as the first element of a basis \(S'\) of \(B'\) , and one would then have \(\mathrm{Tr}_{\mathbf{B}'/\mathbf{A}'}(x'y') = 0\) for all \(y' \in S'\) since a nilpotent endomorphism has all its eigenvalues zero (chap.~VII, § 5, no. 3, cor. 3 of prop. 8), hence zero trace. It would then follow that \(\mathrm{D}_{\mathbf{B}'/\mathbf{A}'}(\mathrm{S}') = 0\) , whence \(\mathrm{D}_{\mathbf{B}'/\mathbf{A}'}(\mathrm{S}) = 0\) (prop. 1) and \(\mathrm{D}_{\mathbf{B}/\mathbf{A}}(\mathrm{S}) = 0\) (formula (6)), contrary to the hypothesis.

Remark. --- Suppose that B is an overfield of A. Let \(S = (x_{1}, \ldots, x_{n})\) be a basis of B, and \((s_{1}, \ldots, s_{n})\) the A-isomorphisms of B into the algebraic closure A' of A (each \(s_{j}\) being repeated \([B : A]_{i}\) times). Recall that, for every \(z \in B\) , one has \(\operatorname{Tr}_{\mathbf{B}/\mathbf{A}}(z) = \sum_{j=1}^{n} s_{j}(z)\) (chap.~VIII, § 12, no. 2, prop. 4).

It then follows from the product formula for determinants that one has

\[
(\det (s _ {j} (x _ {i}))) ^ {2} = \mathrm{D} _ {\mathrm{B/A}} (x _ {1}, \dots , x _ {n}).\tag{7}
\]

This formula shows that proposition 5 generalizes the separability condition given in chap.~V, § 7, no. 2, Remark.

PROPOSITION 6. --- Let \(\Phi\) be an A-bilinear form on E, and K a subring of A such that A is a free K-module of finite dimension q. If \((e_i)_{i=1,\ldots,n}\) is a basis of E over A and \((a_j)_{j=1,\ldots,q}\) a basis of A over K, then \((a_j e_i)\) is a basis of E over K. The map \(\Phi'\) from E \(\times\) E to K defined by \(\Phi'(x,y)=\mathrm{Tr}_{\mathrm{A/K}}(\Phi(x,y))\) is a K-bilinear form on E, and one has

\[
\mathrm{D} _ {\Phi^ {\prime}} (a _ {j} e _ {i}) = \mathrm{N} _ {\mathrm{A} / \mathrm{K}} (\mathrm{D} _ {\Phi} (e _ {1}, \dots , e _ {n})) \cdot (\mathrm{D} _ {\mathrm{A} / \mathrm{K}} (a _ {1}, \dots , a _ {q})) ^ {n}.\tag{8}
\]

The first two assertions being evident, it suffices to prove (7). By definition, the left-hand side is the determinant of the K-linear endomorphism u of E defined by

\[
u (a _ {j} e _ {i}) = \sum_ {r, s} \mathrm{Tr} _ {\mathbf {A} / \mathbf {K}} (a _ {j} a _ {r} \Phi (e _ {i}, e _ {s})) a _ {r} e _ {s}.
\]

Consider the A-linear endomorphism \(\nu\) of E defined by \(\nu(e_i) = \sum_s \Phi(e_i, e_s)e_s\) , and the K-linear endomorphism \(w\) of E defined by \(w(a_j e_i) = (\sum_r \mathrm{Tr}_{\mathrm{A/K}}(a_j a_r)a_r)e_i\) . We have

\[
w \left(v \left(a _ {j} e _ {i}\right)\right) = w \left(\sum_ {s} a _ {j} \Phi \left(e _ {i}, e _ {s}\right) e _ {s}\right) = \sum_ {\tau , s} \operatorname{Tr} _ {\mathrm{A} / \mathrm{K}} \left(a _ {j} \Phi \left(e _ {i}, e _ {s}\right) a _ {\tau}\right) a _ {r} e _ {s}
\]

since \(w(ae_s) = \sum_r \mathrm{Tr}_{\Lambda/\mathbb{K}}(aa_r)a_r e_s\) for all \(a \in A\) ; thus \(u\) is the composite map \(w \circ v\) . Thus, denoting by \(v_{\kappa}\) the map \(v\) considered as a K-linear map, we have \(\det(u) = \det(v_{\kappa})\det(w)\) . Now we have \(\det(v_{\kappa}) = N_{\Lambda/\mathbb{K}}(\det(v))\) (chap.~VIII, § 12, no. 2, prop. 7), and it is clear that \(\det(v) = D_{\Phi}(e_1, \ldots, e_n)\) . On the other hand, since each of the A-modules \(Ae_i\) ( \(i = 1, \ldots, n\) ) is stable under \(w\) and the determinant of the restriction of \(w\) to \(Ae_i\) is \(\det(\mathrm{Tr}_{\Lambda/\mathbb{K}}(a_j a_r)) = D_{\Lambda/\mathbb{K}}(a_1, \ldots, a_q)\) , one has \(\det(w) = (D_{\Lambda/\mathbb{K}}(a_1, \ldots, a_q))^n\) . Formula (8) therefore reduces to \(\det(u) = \det(v_{\kappa})\det(w)\) , a formula proved above.

COROLLARY ("Transitivity formula for discriminants"). — Let K be a commutative ring, A a commutative algebra admitting a finite basis \((a_{j})_{j=1,\ldots,q}\) over K, and E an algebra over A admitting a finite basis \((e_{i})_{i=1,\ldots,n}\) . Then \((a_{j}e_{i})\) is a basis of E over K, and one has

\[
\mathrm{D} _ {\mathrm{E} / \mathrm{K}} (a _ {j} e _ {i}) = \mathrm{N} _ {\mathrm{A} / \mathrm{K}} (\mathrm{D} _ {\mathrm{E} / \mathrm{A}} (e _ {1}, \dots , e _ {n})) \cdot (\mathrm{D} _ {\mathrm{A} / \mathrm{K}} (a _ {1}, \dots , a _ {q})) ^ {n}.
\]

Indeed, if one sets \(\Phi(x,y)=\mathrm{Tr}_{\mathrm{E/A}}(xy)\) , the K-bilinear form \(\Phi'\) of Prop. 6 is \(\Phi'(x,y)=\mathrm{Tr}_{\mathrm{E/K}}(xy)\) by the transitivity formula for traces (Chap. VIII, § 12, no. 2, cor. of Prop. 7).

Exercises. — 1) Let A be an algebra of finite rank over a commutative field K, having an identity element.

\begin{enumerate}
\def\labelenumi{\alph{enumi})}
\item
  Show that if the radical of A is not zero, the bilinear form \((x, y) \to \mathrm{Tr}_{\mathbf{A}/\mathbf{K}}(xy)\) on A is degenerate.
\item
  Assume that K has characteristic 0. Show that if A is a matrix algebra \(\mathbf{M}_{n}(\mathrm{K})\) , S the canonical basis of A over K, we have \(\mathrm{D}_{\mathrm{A}/\mathrm{K}}(\mathrm{S}) \neq 0\) .
\item
  From a) and b), deduce that, for an algebra A of finite rank over a field K of characteristic 0 to be absolutely semisimple, it is necessary and sufficient that the bilinear form \((x,y)\to\mathrm{Tr}_{\mathbf{A}/\mathbf{K}}(xy)\) be nondegenerate (or, equivalently, that \(\mathrm{D}_{\mathbf{A}/\mathbf{K}}(\mathrm{S})\neq0\) for every basis S of A over K).
\end{enumerate}

Let B be a ring, A a subring of B containing the identity element of B; B is therefore an (A, A)-bimodule; we denote by \(^s\text{B}\) (resp. \(^d\text{B}\) ) the set B considered as a left (resp. right) A-module, by \(^s\text{B}^*\) (resp. \(^d\text{B}^*\) ) the right (resp. left) dual A-module of \(^s\text{B}\) (resp. \(^d\text{B}\) ). For every \(x' \in ^s\text{B}^*\) and every \(b \in \text{B}\) , \(x \to \langle xb, x' \rangle\) is an A-linear form on \(^s\text{B}\) , hence an element of \(^s\text{B}^*\) denoted by \(bx'\) ; the map \((b, x') \to bx'\) defines on \(^s\text{B}^*\) a left B-module structure (cf.~Chap.~III, 2nd ed., App. II, no. 7).

\[
(\mathrm{A}, \mathrm{A})
\]

\[
(\mathrm{A}, \mathrm{A}) -
\]

\[
\Phi : (x, y) \rightarrow \varphi (x y)
\]

\(^{s}\) B × \(^{d}\) B into A be nondegenerate, it is necessary and sufficient that \(\varphi(0)\) contain no ideal (left or right) of B other than \(\{0\}\) . We then say that \(\varphi\) is a Frobenius homomorphism from B to A.

\begin{enumerate}
\def\labelenumi{\alph{enumi})}
\setcounter{enumi}{1}
\item
  Let \(\varphi\) a Frobenius homomorphism from B to A; show that the map \(d_{\Phi}\) right-associated to \(\Phi\) is an isomorphism of the left B-module \(B_s\) onto a submodule of the left B-module \(^sB^*\) . Show that \(d_{\Phi}\) is bijective in each of the following two cases: \(1^\circ\) A is a left and right artinian ring satisfying the conditions ( \(N_s\) ) and ( \(N_d\) ) ( \(\S 1\) , exerc. 11), and \(^sB\) and \(^dB\) are free A-modules of finite length (use exerc. 11 b) of \(\S 1\) ); \(2^\circ\) A is a commutative and involutive Artinian ring (chap.~VIII, \(\S 3\) , exerc. 11), contained in the center of B, and \(^{s}\) B is an A-module of finite length (use exerc. 11 of chap.~VIII, § 3).
\item
  Conversely, show that if B \(_{s}\) and \(^{s}\) B \(^{*}\) are isomorphic, there exists a Frobenius homomorphism from B to A when one is in one of the two cases considered in b) and that \(^{s}\) B and \(^{d}\) B have the same length (use exerc. 11 b) of § 1).
\item
  With the hypotheses and notations of a), show that the right annihilator (resp. left annihilator) of a left ideal I (resp. of a right ideal r) of B is the orthogonal l⁰ (resp. r⁰) for the form Φ of the submodule I (resp. r) of sB (resp. dB).
\item
  Let \(\varphi\) a Frobenius homomorphism from B to A. Show that if A is an involutive Artinian ring (chap.~VIII, § 3, exerc. 11) then so is B, when one moreover supposes that one or the other of the conditions of \(b\) ) is satisfied (use \(b\) ) and \(d\) ).
\end{enumerate}

¶ 3) Let A be a commutative field, B an algebra of finite rank over A, having an identity element.

\begin{enumerate}
\def\labelenumi{\alph{enumi})}
\item
  For B to be a Frobenius algebra, it is necessary and sufficient that there exist a Frobenius homomorphism from B to A (exerc. 2). (Use exerc. 2 c) and e) above, and exerc. 6 b) of chap.~VIII, § 13).
\item
  Let \(\varphi\) a Frobenius homomorphism from B to A; every A-linear form on B can then be written in a unique way \(x \to \varphi(b'x)\) (resp. \(x \to \varphi(xb'')\) ) where \(b'\) and \(b''\) belong to B; for this form to be a Frobenius homomorphism, it is necessary and sufficient that \(b'\) (resp. \(b''\) ) be invertible in B.
\item
  For every \(x \in B\) , let \(x^{\sigma}\) be the unique element (cf.~b)) such that \(\varphi(xy) = \varphi(yx^{\sigma})\) for all \(y \in B\) . Show that \(x \to x^{\sigma}\) is an A-automorphism of B. One says that the Frobenius algebra B is symmetric if \(\sigma\) is an inner automorphism of B; there is then a Frobenius homomorphism from B to A for which \(\sigma\) is the identity (cf.~b)). This is equivalent to saying that the (B, B)-bimodules B and \(^{s}B^{*} = ^{d}B^{*}\) (written B*) are isomorphic (exerc. 2 c)).
\item
  Let E be a left B-module of finite length, E' its dual, E'* the dual of E' considered as a vector space over A; E'* is endowed with a structure of left B-module by setting, for \(x' \in \mathrm{E}', x'' \in \mathrm{E}*\) , \(b \in \mathrm{B}\) , \(\langle x', bx'' \rangle = \langle x'b, x'' \rangle\) (chap.~III, 2nd ed., App. II, no. 7). For every \(x \in \mathrm{E}\) , let \(f_{\mathrm{E}}(x)\) (or simply \(f(x)\) ) be the element of E'* such that \(\langle x', f(x) \rangle = \varphi(\langle x, x' \rangle)\) for all \(x' \in \mathrm{E}'\) ; show that \(f\) is a semilinear bijection for the automorphism \(\sigma\) , from the left B-module E onto the left B-module E'* (use exerc. 10 of chap.~VIII, § 4). For \(\mathrm{E} = \mathrm{B}_s\) , one has (with the notations of exerc. 2 b)) \(d_{\Phi}(x^{\sigma}) = f_{\mathrm{B}_s}(x)\) for all \(x \in \mathrm{B}\) .
\end{enumerate}

\begin{enumerate}
\def\labelenumi{\arabic{enumi})}
\setcounter{enumi}{3}
\tightlist
\item
  \begin{enumerate}
  \def\labelenumii{\alph{enumii})}
  \tightlist
  \item
    Let G be a finite group. Show that the algebra B of the group G over any commutative field A is a symmetric Frobenius algebra (exerc. 3). (Consider the map \(\varphi\) from B to A which, to every element \(x = \sum_{s \in G} \xi_s.s\) , associates \(\varphi(x) = \xi_e\) , \(e\) denoting the identity element of G).
  \end{enumerate}
\end{enumerate}

\begin{enumerate}
\def\labelenumi{\alph{enumi})}
\setcounter{enumi}{1}
\tightlist
\item
  Let E be a vector space of finite dimension n over a commutative field A, and B the exterior algebra \(\wedge\) E. Show that B is a Frobenius algebra. (If \((e_i)_{1 \leqslant i \leqslant n}\) is a basis of E, consider the map which to every element \(x\) of B associates the coefficient of \(e_1 \wedge e_2 \wedge \ldots \wedge e_n\) in the expression of \(x\) using the basis of B corresponding to \((e_i)\) ). For the Frobenius algebra B to be symmetric, it is necessary and sufficient that \(n\) be even or A have characteristic 2.
\end{enumerate}

¶ 5) a) Show that the tensor product of two Frobenius algebras (resp. Frobenius and symmetric algebras) of finite rank over a commutative field K is a Frobenius algebra (resp. Frobenius and symmetric algebra) (cf.~§ 1, no. 9, prop. 9).

\begin{enumerate}
\def\labelenumi{\alph{enumi})}
\setcounter{enumi}{1}
\tightlist
\item
  Let B be an algebra of finite rank over K, L an extension of K of finite degree over K. Show that if the algebra \(\mathrm{B}_{(\mathrm{L})} = \mathrm{B} \otimes_{\mathrm{K}} \mathrm{L}\) over L is Frobenius (resp. Frobenius and symmetric), then so is B. (Use exerc. 2 c) and 3 c) above and exerc. 2 of chap.~VIII, § 2).
\end{enumerate}

\begin{enumerate}
\def\labelenumi{\arabic{enumi})}
\setcounter{enumi}{5}
\tightlist
\item
  Show that every absolutely semisimple algebra B of finite rank over a commutative field A is a symmetric Frobenius algebra. (Reduce to the case where B is simple; use exerc. 5 b) above, as well as prop. 9 of chap.~VIII, § 12, no. 3).
\end{enumerate}

\section{Hermitian forms and quadratic forms}

In all that follows in this Chapter, unless expressly stated otherwise, A denotes a ring and E a left A-module. One supposes A endowed with an involutive antiautomorphism J, denoted \(\alpha \to \overline{\alpha}\) ; one therefore has \((\overline{\alpha + \beta}) = \overline{\alpha} + \overline{\beta}\) , \((\overline{\alpha\beta}) = \overline{\beta}.\overline{\alpha}\) and \(\overline{\overline{\alpha}} = \alpha\) for all \(\alpha\) , \(\beta\) in A. Unless expressly stated otherwise, the sesquilinear forms considered are right sesquilinear ( \(\S 1\) , no. 2, def. 4) for this antiautomorphism.

\subsection{\texorpdfstring{Hermitian and \(\varepsilon\)-hermitian forms.}{Hermitian and epsilon-hermitian forms.}}

DEFINITION 1. --- Let \(\varepsilon\) be an element of the center of A. A sesquilinear form \(\Phi\) on E such that one has \(\Phi(x, y) = \varepsilon \overline{\Phi(y, x)}\) for all \(x\) and \(y\) in E is called a form \(\varepsilon\) -hermitian. A 1-hermitian form (resp. \((-1)\) -hermitian) is said to be hermitian (resp. antihermitian).

When J is the identity (which implies that A is commutative), a Hermitian (resp. antihermitian) form (for J) is none other than a symmetric (resp. antisymmetric) bilinear form (chap.~III, § 5, no. 1, def. 2). Recall that an alternating bilinear form (chap.~III, § 5, no. 2, def. 4) is antisymmetric; the converse is true if, in A, the relation \(2a = 0\) implies \(a = 0\) .

The orthogonality relation (§ 1, no. 3) with respect to an ε-hermitian form is evidently symmetric (cf.~exerc. 1).

If \(\alpha\) is an invertible element of A, the map T: \(\lambda \to \alpha^{-1}\overline{\lambda}\alpha\) is an antiautomorphism of A, and one readily verifies that the form \(\Phi\alpha\) is sesquilinear with respect to T. If, moreover, we have \(\alpha = \overline{\alpha}\) then T is involutive, and, if \(\Phi\) is \(\varepsilon\) -hermitian, so is \(\Phi\alpha\) ; indeed we have

\[
(\lambda^ {\mathrm{T}}) ^ {\mathrm{T}} = \alpha^ {- 1} (\overline {{\alpha^ {- 1} \lambda \alpha}}) \alpha = \alpha^ {- 1} \overline {{\alpha}} \lambda \overline {{\alpha}} ^ {- 1} \alpha = \lambda
\]

\[
\Phi (y, x) \alpha = \varepsilon \overline {{\Phi (x , y)}} \alpha = \varepsilon (\Phi (x, y) \alpha) ^ {\mathrm{T}}.
\]

In particular, when A is a field, the elements \(\alpha\) of the center of A such that \(\bar{\alpha} = \alpha\) form a subfield K of A, and the \(\varepsilon\) -hermitian forms on E (for J) form a vector space over K.

Remarks. --- 1) If \(\Phi\) is a form \(\varepsilon\) -hermitian on E, we have \(\Phi(x, y) = \varepsilon \Phi(x, y) \bar{\varepsilon}\) for all \(x, y\) in E. Therefore, if \(\Phi\) takes invertible values, we have \(\varepsilon \bar{\varepsilon} = 1\) .

\begin{enumerate}
\def\labelenumi{\arabic{enumi})}
\setcounter{enumi}{1}
\tightlist
\item
  If there exists an invertible element i of the center of A such that \(\bar{i}=i\varepsilon\) , then, for \(\Phi\) to be \(\varepsilon\) -hermitian, it is necessary and sufficient that \(i\Phi\) be Hermitian.
\end{enumerate}

The map \((y, x) \to \overline{\Phi(x, y)}\) being sesquilinear for J, for \(\Phi\) to be \(\varepsilon\) -hermitian, it is necessary and sufficient that one has \(\Phi(y, x) = \varepsilon \overline{\Phi(x, y)}\) when \(x\) and \(y\) range over a system of generators of E. In particular, if E admits a finite basis \((e_i)_{1 \leqslant i \leqslant n}\) for a sesquilinear form \(\Phi\) on E to be \(\varepsilon\) -hermitian, it is necessary and sufficient that its matrix \(R = (\rho_{ij}) = (\Phi(e_i, e_j))\) satisfies the relations \(\rho_{ji} = \varepsilon \overline{\rho_{ij}}\) for all \(i, j\) , that is to say \(^t R = \varepsilon \overline{R}\) ; a matrix \(R\) possessing this property is said to be \(\varepsilon\) -Hermitian. When \(\varepsilon = 1\) (resp. -1) one says that \(R\) is Hermitian (resp. antihermitian) with respect to the anti-automorphism J. When J is the identity (hence A is commutative), a Hermitian (resp. antihermitian) matrix \(R\) is such that \(^t R = R\) (resp. \(^t R = -R\) ); we then say that \(R\) is a symmetric (resp. antisymmetric) matrix. For \(\Phi\) to be an alternating form, it is necessary and sufficient that its matrix be antisymmetric and, in addition, that the diagonal terms of R all be zero; a matrix possessing these properties is said to be alternating.

Let \(\Phi\) a sesquilinear form on E, and let \(s_{\Phi}\) and \(d_{\Phi}\) be the maps from E into E* associated with \(\Phi\) to the left and to the right (§ 1, no. 6). For \(\Phi\) to be \(\varepsilon\) -hermitian, it is necessary and sufficient that \(\langle x, s_{\Phi}(y) \rangle = \bar{\varepsilon} \langle x, d_{\Phi}(y) \rangle\) for all elements \(x, y\) of E, hence that \(s_{\Phi} = \bar{\varepsilon} d_{\Phi}\) , or equivalently that \(\langle x, d_{\Phi}(y) \rangle = \varepsilon \langle x, s_{\Phi}(y) \rangle\) , hence that \(d_{\Phi} = \varepsilon s_{\Phi}\) .

Let \(\Phi\) be a form \(\varepsilon\) -hermitian such that the map \(d_{\Phi}\) from E into E* associated on the right to \(\Phi\) is bijective. For every endomorphism \(u\) of E we then have

\[
u ^ {* *} = \varepsilon \bar {\varepsilon} u.\tag{1}
\]

Indeed, whatever the elements \(x\) and \(y\) of \(\mathbf{E}\) , one has

\[
\begin{array}{r l} \Phi (x, u ^ {* *} (y)) & = \Phi (u ^ {*} (x), y) = \varepsilon \overline {{\Phi (y , u ^ {*} (x))}} = \varepsilon \overline {{\Phi (u (y) , x)}} \\ & = \varepsilon \Phi (x, u (y)) \bar {\varepsilon} = \Phi (x, \varepsilon \bar {\varepsilon} u (y)) \end{array}
\]

therefore \(u^{**}(x) = \varepsilon \bar{\varepsilon} u(x)\) since \(\Phi\) is nondegenerate.

If \(\Phi\) is a form \(\varepsilon\) -hermitian such that the maps \(s_{\Phi}\) and \(d_{\Phi}\) are bijective, then the inverse form \(\hat{\Phi}\) of \(\Phi\) ( \(\S 1\) , no. 7) is a \(\bar{\varepsilon}\) -hermitian form. Indeed, setting \(s = s_{\Phi}\) , \(d = d_{\Phi}\) for brevity, one deduces from \(d = \varepsilon s\) that \(s^{-1} = \bar{\varepsilon} d^{-1}\) , \(s\) being semilinear. Consequently, whatever \(u\) , \(\nu\) in E, one has

\[
\hat {\Phi} (u, \nu) = \Phi (s ^ {- 1} (u), d ^ {- 1} (\nu)) = \bar {\varepsilon} \Phi (d ^ {- 1} (u), d ^ {- 1} (\nu)),
\]

whence

\[
\widehat {\Phi} (\nu , u) = \overline {{\varepsilon \varepsilon}} \overline {{\Phi (d ^ {- 1} (u) , d ^ {- 1} (\nu))}} = \overline {{\varepsilon}} \widehat {\Phi} (u, \nu),
\]

since ε is in the center of A.

Finally, when the ring A is commutative, the canonical extensions of a \(\varepsilon\) -hermitian form \(\Phi\) to the tensor and exterior powers \(\bigotimes^{p}\mathbf{E}\) and \(\bigwedge^{p}\mathbf{E}\) of E are \(\varepsilon^{p}\) -hermitian forms, as follows immediately from formulas (35) and (37) of § 1, no. 9.

\subsection{Modules over a quadratic extension.}

Let K be a commutative ring. One takes for A the quadratic extension \(\mathrm{A} = \mathrm{K}(i)\) with \(i^{2} = -1\) , and for J the automorphism \(\lambda + i\mu \to \lambda - i\mu\) ( \(\lambda \in \mathbf{K}\) , \(\mu \in \mathbf{K}\) ) (chap.~II, § 7, no. 7). If E is an A-module, we denote by \(\mathbf{E}_0\) the K-module deduced from E by restriction of the ring of scalars, and by \(j\) the automorphism \(x \to ix\) of \(\mathbf{E}_0\) ; one evidently has \(j^2 = -I\) , where I is the identity map of \(\mathbf{E}_0\) . Conversely, let \(\mathbf{E}_0\) be a K-module and let \(j\) be an automorphism of \(\mathbf{E}_0\) such that \(j^2 = -I\) ; the map \(\lambda + i\mu \to \lambda I + \mu j\) is evidently a homomorphism from A into the ring \(\mathfrak{L}(\mathbf{E}_0)\) of endomorphisms of \(\mathbf{E}_0\) ; we have therefore defined on \(\mathbf{E}_0\) a structure of A-module, for which we have

\[
(\lambda + i \mu) x = \lambda x + \mu j (x) \quad (x \in \mathrm{E} _ {0}, \lambda \in \mathrm{K}, \mu \in \mathrm{K}).\tag{2}
\]

If E' is another A-module, \(E_{0}'\) the underlying K-module of \(E'\) , \(j'\) the automorphism \(x' \rightarrow ix'\) of \(E_{0}'\) then the A-linear maps f from E into \(E'\) are none other than the K-linear maps from \(E_{0}\) into \(E_{0}'\) such that \(f \circ j = j' \circ f\) . In particular, if one denotes by \(E^{*}\) and \((\mathrm{E}_{0})^{*}\) the respective duals of E and \(E_{0}\) , and if \(f_{1}\) and \(f_{2}\) are two maps from E into K, for the map \(x \to f_{1}(x) + if_{2}(x)\) from E into A to be A-linear, it is necessary and sufficient that \(f_{1}\) and \(f_{2}\) be in \((\mathrm{E}_{0})^{*}\) and that one has \(f_{1} \circ j + i(f_{2} \circ j) = if_{1} - f_{2}\) , that is to say \(f_{1} = f_{2} \circ j\) and \(f_{1} \circ j = -f_{2}\) . Since j is an automorphism of \(E_{0}\) and \(j^{2} = -I\) , these two conditions are equivalent. Eliminating \(f_{1}\) or \(f_{2}\) , we see that the formulas

\begin{enumerate}
\def\labelenumi{(\arabic{enumi})}
\setcounter{enumi}{2}
\tightlist
\item
\end{enumerate}

\[
f (x) = f _ {1} (x) - i f _ {1} (j (x))\tag{4}
\]

\[
f (x) = f _ {2} (j (x)) + i f _ {2} (x)
\]

\((x \in \mathbf{E}, f \in \mathbf{E}^{*}, f_{1} \in (\mathbf{E}_{0})^{*}, f_{2} \in (\mathbf{E}_{0})^{*})\) establish two bijective correspondences between \(\mathbf{E}^{*}\) and \((\mathbf{E}_{0})^{*}\) .

\subsection{Bilinear forms associated with a Hermitian form.}

We assume again here that the ring A is the quadratic extension \(\mathrm{A} = \mathrm{K}(i)\) (where \(i^{2} = -1\) ) of a commutative ring K, and that J is the automorphism \(\lambda + i\mu \to \lambda - i\mu\) of A ( \(\lambda \in K, \mu \in K\) ). Let E and \(E'\) be two A-modules, \(E_{0}\) and \(E_{0}'\) be the underlying K-modules of E and \(E'\) , j and \(j'\) the automorphisms \(x \to ix\) and \(x' \to ix'\) of E and \(E'\) (cf.~no. 2). A K-bilinear form f on \(E_{0} \times E_{0}'\) will be said to be invariant under j and \(j'\) if one has

\[
f (j (x), j ^ {\prime} (x ^ {\prime})) = f (x, x ^ {\prime})\tag{5}
\]

for \(x \in \mathrm{E}_0\) and \(x' \in \mathrm{E}_0'\) . Replacing \(x\) by \(j(x)\) , one sees that this condition is equivalent to

\[
f (x, j ^ {\prime} (x ^ {\prime})) = - f (j (x), x ^ {\prime})\tag{6}
\]

for all \(x \in \mathbf{E}_0\) and \(x' \in \mathbf{E}_0'\) .

PROPOSITION 1. --- Let \(\Phi_{1}\) (resp. \(\Phi_{2}\) ) be a K-bilinear form on \(\mathbf{E}_{0} \times \mathbf{E}_{0}^{\prime}\) , invariant under \(j\) and \(j^{\prime}\) . The map which associates to \(\Phi_{1}\) (resp. \(\Phi_{2}\) ) the map \(\Phi\) of \(\mathbf{E} \times \mathbf{E}^{\prime}\) into A defined by

\begin{enumerate}
\def\labelenumi{(\arabic{enumi})}
\setcounter{enumi}{6}
\tightlist
\item
\end{enumerate}

\[
\Phi (x, x ^ {\prime}) = \Phi_ {1} (x, x ^ {\prime}) + i \Phi_ {1} (x, j ^ {\prime} (x ^ {\prime}))\tag{8}
\]

\[
(\text { resp. } \Phi (x, x ^ {\prime}) = - \Phi_ {2} (x, j ^ {\prime} (x ^ {\prime})) + i \Phi_ {2} (x, x ^ {\prime}))
\]

\((x \in \mathrm{E}, x' \in \mathrm{E})\) , is an isomorphism of the K-vector space of K-bilinear forms on \(\mathrm{E}_0 \times \mathrm{E}_0'\) invariant under \(j\) and \(j'\) onto the K-vector space of sesquilinear forms on \(\mathrm{E} \times \mathrm{E}'\) . Suppose moreover \(\mathrm{E} = \mathrm{E}'\) ; for \(\Phi\) to be Hermitian, it is necessary and sufficient that \(\Phi_1\) be symmetric (resp. that \(\Phi_2\) be antisymmetric) (cf.~exerc. 4).

Indeed every map \(\Phi\) of E × E' into A can be written, in one and only one way, in the form \(\Phi = \Phi_1 + i\Phi_2\) , where \(\Phi_1\) and \(\Phi_2\) are maps of E × E' into K. For the partial map \(x \to \Phi(x, x')\) to be A-linear, it is necessary and sufficient, according to formula (3) (resp. (4)) of no. 2, that \(\Phi_1\) (resp. \(\Phi_2\) ) be K-linear in \(x\) and that one has

\begin{enumerate}
\def\labelenumi{(\arabic{enumi})}
\setcounter{enumi}{8}
\tightlist
\item
\end{enumerate}

\[
\Phi (x, x ^ {\prime}) = \Phi_ {1} (x, x ^ {\prime}) - i \Phi_ {1} (j (x), x ^ {\prime})\tag{10}
\]

\[
(\text { resp. } \Phi (x, x ^ {\prime}) = \Phi_ {2} (j (x), x ^ {\prime}) + i \Phi_ {2} (x, x ^ {\prime})).
\]

Similarly, for \(\overline{\Phi(x,x^{\prime})}\) to be A-linear in \(x^{\prime}\) , it is necessary and sufficient that \(\Phi_{1}\) (resp. \(\Phi_{2}\) ) be K-linear in \(x^{\prime}\) and that one has

\begin{enumerate}
\def\labelenumi{(\arabic{enumi})}
\setcounter{enumi}{10}
\tightlist
\item
\end{enumerate}

\[
\Phi (x, x ^ {\prime}) = \Phi_ {1} (x, x ^ {\prime}) + i \Phi_ {1} (x, j ^ {\prime} (x ^ {\prime}))\tag{12}
\]

\[
(\text { resp. } \Phi (x, x ^ {\prime}) = - \Phi_ {2} (x, j ^ {\prime} (x ^ {\prime})) + i \Phi_ {2} (x, x ^ {\prime})).
\]

It follows immediately that, for \(\Phi\) to be sesquilinear, it is necessary and sufficient that it be written in one or the other of the forms (9) and (11) (resp. (10) and (12)) with \(\Phi_{1}\) (resp. \(\Phi_{2}\) ) K-bilinear invariant under \(j\) and \(j'\) .

It follows from this that, for a sesquilinear form \(\Phi = \Phi_{1} + i\Phi_{2}\) to be zero, it is necessary and sufficient that \(\Phi_{1}\) (resp. \(\Phi_{2}\) ) be zero. Now, if \(E = E'\) , one has \(\Phi(y, x) = \Phi_{1}(y, x) + i\Phi_{2}(y, x)\) and \(\overline{\Phi(x, y)} = \Phi_{1}(x, y) - i\Phi_{2}(x, y)\) for these two expressions to be equal, in other words for \(\Phi\) to be Hermitian, it is therefore necessary and sufficient that \(\Phi_{1}\) be symmetric (resp. that \(\Phi_{2}\) be antisymmetric).

Remarks. --- 1) Formulas (7) and (8) show that, if \(x \in E\) , for \(\Phi(x, x') = 0\) to hold for all \(x' \in E'\) , it is necessary and sufficient that \(\Phi_1(x, x') = 0\) (resp. \(\Phi_2(x, x') = 0\) ) for all \(x' \in E'\) .

\begin{enumerate}
\def\labelenumi{\arabic{enumi})}
\setcounter{enumi}{1}
\tightlist
\item
  The adjoint of an endomorphism \(u\) of E with respect to \(\Phi\) ( \(\S 1\) , no. 8) is the same as the adjoint of \(u\) (considered as an endomorphism of \(\mathrm{E}_{0}\) ) with respect to \(\Phi_{1}\) (resp. \(\Phi_{2}\) ).
\end{enumerate}

\subsection{Quadratic forms.}

DEFINITION 2. --- Suppose the ring A is commutative. A map Q from E to A is said to be a quadratic form on E if

\begin{enumerate}
\def\labelenumi{\arabic{enumi})}
\item
  one has Q(αx) = α²Q(x) for α ∈ A and x ∈ E ;
\item
  the map \(\Phi : (x, y) \to \mathrm{Q}(x + y) - \mathrm{Q}(x) - \mathrm{Q}(y)\) of \(\mathbf{E} \times \mathbf{E}\) in A is a bilinear form.
\end{enumerate}

The bilinear form \(\Phi\) (which is necessarily symmetric) is called the bilinear form associated with Q. If \(\Phi\) is nondegenerate, Q is said to be nondegenerate.

Since Q(2x) = 4Q(x), it follows immediately from 2) that we have

\[
\Phi (x, x) = 2 \mathrm{Q} (x).\tag{13}
\]

In particular, if A is a ring of characteristic 2, the form \(\Phi\) is alternating.

We will say that two elements (resp. two subsets) of E are orthogonal with respect to Q if they are orthogonal with respect to the associated bilinear form Φ.

Let \((x_{i})_{i\in I}\) be a family of elements of E and \((a_{i})_{i\in I}\) a family of elements of A, all but finitely many of which are zero. By induction on the number of indices \(i\) for which \(a_{i}\neq0\) one shows easily that we have

\[
\mathrm{Q} (\sum_ {i} a _ {i} x _ {i}) = \sum_ {i} a _ {i} ^ {2} \mathrm{Q} (x _ {i}) + \sum_ {\{i, j \}} a _ {i} a _ {j} \Phi (x _ {i}, x _ {j}),\tag{14}
\]

the last summation being extended to the two-element subsets of I.

For every bilinear form \(f\) on \(\mathbf{E} \times \mathbf{E}\) one defines a quadratic form \(\mathbf{Q}\) by setting \(\mathrm{Q}(x) = f(x, x)\) ; the bilinear form \(\Phi\) associated with \(\mathbf{Q}\) is then defined by \(\Phi(x, y) = f(x, y) + f(y, x)\) for \(x, y\) in \(\mathbf{E}\) . Moreover, if one assumes that the scalar 2 has an inverse \(\frac{1}{2}\) in \(\mathbf{A}\) , there exists one and only one symmetric bilinear form \(f\) such that \(\mathrm{Q}(x) = f(x, x)\) , namely \(f = \frac{1}{2}\Phi\) ; the discriminant of \(f\) with respect to a system \(S = (x_1, \ldots, x_n)\) is also called the discriminant of \(\mathbf{Q}\) with respect to \(S\) . In this case, there is therefore a one-to-one correspondence between quadratic forms and symmetric bilinear forms on \(\mathbf{E}\) (cf. exercise 6).

In the case of a free module, we also have the following result:

PROPOSITION 2. --- Suppose that A is commutative and that E admits a basis \((e_i)_{i \in I}\) . Then, for every quadratic form Q on E, there exists a bilinear form \(f\) on E \(\times\) E such that Q(x) = f(x, x) for all \(x \in E\) . For every family \((b_{ij})_{(i,j) \in I \times I}\) of scalars such that \(b_{ij} = b_{ji}\) for \((i, j) \in I \times I\) , there exists one and only one quadratic form Q such that

\[
\mathrm{Q} (e _ {i}) = b _ {i i}, \quad \Phi (e _ {i}, e _ {j}) = b _ {i j} \text {   for   } i \neq j,\tag{15}
\]

where Φ denotes the bilinear form associated with Q; then Q is given by the formula

\[
\mathrm{Q} (\sum_ {i} a _ {i} e _ {i}) = \sum_ {\{i, j \}} b _ {i j} a _ {i} a _ {j},\tag{16}
\]

the last summation being extended over the subsets \(\{i, j\}\) of I having one or two elements.

Indeed, since formula (16) is only a transcription of formula (14), the uniqueness of a quadratic form Q satisfying (15) is proved. To prove its existence, we first note that there exists a family ( \(b_{ij}'\) ) of elements of A such that \(b_{ii}' = b_{ii}\) and \(b_{ij}' + b_{ji}' = b_{ij}\) for \(i \neq j\) ; for example, one obtains such a family by endowing I with a structure of totally ordered set (Ens., chap.~III, § 2, no. 3, th. 1) and setting \(b_{ij}' = b_{ij}\) for \(i<j\) and \(b_{ij}^{\prime}=0\) for \(i>j\) . Since the \(e_{i}\) form a basis of E, there exists a bilinear form \(f\) on \(\mathbf{E}\times\mathbf{E}\) such that \(f(e_{i},e_{j})=b_{ij}^{\prime}\) ; by setting \(\mathrm{Q}^{\prime}(x)=f(x,x)\) and denoting by \(\Phi^{\prime}\) the bilinear form associated with the quadratic form \(\mathrm{Q}^{\prime}\) , we obtain \(\mathrm{Q}^{\prime}(e_{i})=b_{ii}\) and \(\Phi^{\prime}(e_{i},e_{j})=f(e_{i},e_{j})+f(e_{j},e_{i})=b_{ij}\) . This proves our second assertion. As for the first, it follows immediately, because, by virtue of uniqueness, if a quadratic form Q satisfies (15), we have \(\mathrm{Q}(x)=\mathrm{Q}^{\prime}(x)=f(x,x)\) .

The module E equipped with the structure defined by a quadratic form Q is called a quadratic module. A homomorphism of the quadratic module (E, Q) into a quadratic module (E', Q') is a linear map u from E into E' such that Q = Q' ∘ u; if Φ and Φ' are the bilinear forms associated with Q and Q', we then have Φ(x, y) = Φ'(u(x), u(y)) for x ∈ E, y ∈ E; in other words, Φ' is the inverse image of Φ under u (§ 1, no. 1). Two quadratic forms Q and Q' on two A-modules E and E' are said to be equivalent if the corresponding quadratic modules are isomorphic.

Let \((\mathbf{E}_{i}, \mathbf{Q}_{i})_{i\in\mathbb{I}}\) be a family of quadratic modules, and let E be the direct sum of the modules \(E_{i}\) . One calls the external direct sum of the quadratic modules \((\mathbf{E}_{i}, \mathbf{Q}_{i})\) the quadratic module obtained by equipping E with the quadratic form Q defined by \(\mathrm{Q}(\sum_{i}x_{i}) = \sum_{i}\mathrm{Q}_{i}(x_{i})\) for \(x_{i} \in E_{i}\) . One also says that the quadratic form Q is the external direct sum of the quadratic forms \(Q_{i}\) .

If the forms \(Q_{i}\) are nondegenerate, so is Q.

Let Q be a quadratic form on the A-module E; if F is a submodule of E and if Q is constant on each class modulo F, the map \(\overline{Q}\) from E/F into A deduced from Q by passing to the quotient is evidently a quadratic form, and the canonical map from E onto E/F is a homomorphism for the structures of quadratic modules. For Q to be constant on each class modulo F, it is necessary and sufficient that one have \(Q(x + y) = Q(x)\) for \(x \in E\) and \(y \in F\) , that is, denoting by \(\Phi\) the bilinear form associated with Q, that one have \(Q(y) + \Phi(x, y) = 0\) for \(y \in F\) and \(x \in E\) . Setting x = 0, one sees that one has \(Q(y) = 0\) for \(y \in F\) , and hence \(\Phi(x,y)=0\) for \(x\in E\) and \(y\in F\) . In other words, if one calls the kernel of the quadratic module (E, Q) the set N of elements x of E such that \(Q(x)=0\) and \(\Phi(x,z)=0\) for all \(z\in E\) , for Q to be constant on each class modulo F, it is necessary and sufficient that F be contained in the kernel N of (E, Q). One verifies without difficulty that N is a submodule of E. For Q to be constant on each class modulo F, it is therefore necessary and sufficient that F be generated by elements of N.

One sees immediately that the kernel of the quadratic module E/N is \{0\}.

PROPOSITION 3. --- Let h be a homomorphism of A into a commutative ring A'. For every quadratic form Q on the A-module E, there exists one and only one quadratic form Q' on the A'-module A'⊗A E (chap.~III, 2nd ed., App. II, no. 10) such that one has

\[
\mathrm{Q} ^ {\prime} (1 \otimes x) = h (\mathrm{Q} (x))\tag{17}
\]

for all \(x \in \mathbf{E}\) . Moreover, the bilinear form \(\Phi'\) associated with Q' is obtained by extension of the ring of scalars from the bilinear form \(\Phi\) associated with Q.

We first show that, if there exists a quadratic form Q' satisfying (17), it is unique and the bilinear form \(\Phi'\) associated with Q' is obtained by extension of the ring of scalars from the form \(\Phi\) associated with Q. Indeed, this latter assertion follows from the fact that one has

\[
\begin{array}{r l} \Phi^ {\prime} (1 \otimes x, 1 \otimes y) & = Q ^ {\prime} (1 \otimes x + 1 \otimes y) - Q ^ {\prime} (1 \otimes x) - Q ^ {\prime} (1 \otimes y) \\ & = h (\Phi (x, y)) \end{array}\tag{18}
\]

for \(x \in \mathrm{E}, y \in \mathrm{E}\) . Formula (14) then shows that one has

\[
\mathrm{Q} ^ {\prime} \left(\sum_ {i} a _ {i} ^ {\prime} \otimes x _ {i}\right) = \sum_ {i} a _ {i} ^ {\prime 2} h (\mathrm{Q} (x _ {i})) + \sum_ {\{i, j \}} a _ {i} ^ {\prime} a _ {j} ^ {\prime} h (\Phi (x _ {i}, x _ {j}))\tag{19}
\]

for \(a_{i}^{\prime}\in\mathbf{A}^{\prime}\) and \(x_{i}\in\mathbf{E}\) , which proves the uniqueness of \(\mathbf{Q}^{\prime}\) .

To show the existence of Q', we shall first suppose that the module E admits a basis \((e_{i})_{i\in I}\) . There then exist elements \(b_{ij}\) of A such that \(b_{ij}=b_{ji}\) and \(\mathrm{Q}(\sum_{i}a_{i}e_{i})=\sum_{\{i,j\}}b_{ij}a_{i}a_{j}\) for \(a_{i}\in A\) (prop. 2). Since the elements \(1\otimes_{A}e_{i}\) form a basis of the

A'-module A' ⊗\_A E, one defines a quadratic form Q' on this latter module by setting

\[
\mathrm{Q} ^ {\prime} (\sum_ {i} a _ {i} ^ {\prime} \otimes e _ {i}) = \sum_ {\{i, j \}} a _ {i} ^ {\prime} a _ {j} ^ {\prime} h (b _ {i j})\tag{20}
\]

for \(a_i' \in A'\) ; whence, for \(x = \sum_{i} a_i e_i \in E\)

\[
Q ^ {\prime} (1 \otimes x) = Q ^ {\prime} (\sum_ {i} h (a _ {i}) \otimes e _ {i}) = \sum_ {\{i, j \}} h (a _ {i}) h (a _ {j}) h (b _ {i j}) = h (Q (x)),\tag{21}
\]

which proves the existence of Q' in this case.

Now let us turn to the general case. Let \((x_{i})_{i\in I}\) a system of generators of E, \(\mathbf{A}^{(1)}\) the module of formal linear combinations of elements of I (chap.~II, § 1, no. 8), and \((e_{i})_{i\in I}\) the canonical basis of \(\mathbf{A}^{(1)}\) . The linear map \(u\) of \(\mathbf{A}^{(1)}\) in E defined by \(u(e_{i}) = x_{i}\) is surjective since the elements \(x_{i}\) generate E. It follows (chap.~III, 2nd ed., App. II, no. 5, prop. 4) that the map \(1\otimes u\) of \(\mathbf{A}'\otimes\mathbf{A}^{(1)}\) in \(\mathbf{A}'\otimes\mathbf{E}\) is surjective, and its kernel P' is generated by elements of the form \(1\otimes p\) with \(u(p)=0\) Let then \(\mathrm{Q}_{1}^{\prime}\) be the extension to \(\mathbf{A}'\otimes_{\mathbf{A}}\mathbf{A}^{(1)}\) of the quadratic form \(\mathrm{Q}_{1}=\mathrm{Q}\circ u\) on \(\mathbf{A}^{(1)}\) . If \(p\) is an element of \(\mathbf{A}^{(1)}\) such that \(u(p)=0\) , one has \(\mathrm{Q}_{1}^{\prime}(1\otimes p)=h(\mathrm{Q}_{1}(p))=0\) and (denoting by \(\Phi_{1}^{\prime}\) the bilinear form associated with \(\mathrm{Q}_{1}^{\prime}\) ) \(\Phi_{1}^{\prime}(1\otimes p,1\otimes x)=h(\Phi(u(p),u(x)))=0\) for all \(x\in\mathbf{A}^{(1)}\) . Therefore, if \(u(p)=0\) ( \(p\in\mathbf{A}^{(1)}\) ), then \(1\otimes p\) belongs to the kernel of the quadratic module \(\mathbf{A}'\otimes_{\mathbf{A}}\mathbf{A}^{(1)}\) and consequently, as we saw above, there exists a quadratic form Q' on \(\mathbf{A}'\otimes_{\mathbf{A}}\mathbf{E}\) such that \(\mathrm{Q}_{1}^{\prime}=\mathrm{Q}^{\prime}\circ(1\otimes u)\) . Since \(u\) is surjective, we see that Q' satisfies condition (17). QED.

The quadratic form Q′, whose existence and uniqueness are guaranteed by Prop. 3, is called the extension of Q to A′ ⊗\_A E (with respect to h). One also says that Q' is obtained from Q by extension of the ring of scalars.

Exercises. --- 1) Let A be a field, E a left vector space over A, \(\Phi\) a sesquilinear form on E (for an antiautomorphism J of A). We assume that the rank (finite or infinite) of \(\Phi\) is \(\geqslant 2\) and that the relations \(\Phi(x, y) = 0\) and \(\Phi(y, x) = 0\) are equivalent.

\begin{enumerate}
\def\labelenumi{\alph{enumi})}
\item
  Show that there exists \(\lambda\neq0\) in A such that one has \(\Phi(y,x)=\lambda(\Phi(x,y))^{\mathfrak{J}}\) (Use exerc. 8 of § 1).
\item
  Show that there exists \(\alpha\in A\) such that the sesquilinear form \(\Phi\alpha\) (for the antiautomorphism \(\xi\to\alpha^{-1}\xi^{J}\alpha\) ) is Hermitian or alternating.
\end{enumerate}

(First note that one has \(\xi^{J^2} = \lambda^{-1}\xi\lambda^{-J}\) and \(\lambda\lambda^J = \lambda^J\lambda = 1\) . Then distinguish two cases according as \(\xi + \xi^J\lambda^{-1} = 0\) for all \(\xi \in A\) or not; in the second case, show that every element \(\neq 0\) of the form \(\alpha = \xi + \xi^J\lambda^{-1}\) answers the question).

\begin{enumerate}
\def\labelenumi{\arabic{enumi})}
\setcounter{enumi}{1}
\tightlist
\item
  Let \(\Phi\) be a sesquilinear form \(\varepsilon\) -hermitian on a finite-dimensional vector space E over a field A.
\end{enumerate}

\begin{enumerate}
\def\labelenumi{\alph{enumi})}
\item
  Show that for every vector subspace M of E, we have \((\mathbf{M}^{\mathbf{0}})^{\mathbf{0}} = \mathbf{M} + \mathbf{E}^{\mathbf{0}}\) and \(\dim M^{\mathbf{0}} + \dim M = \dim E + \dim (M \cap E^{\mathbf{0}})\) .
\item
  If \(M_{1}\) , \(M_{2}\) are two vector subspaces of E, show that one has dim \((\mathbf{M}_{1} \cap \mathbf{M}_{2}^{0}) + \dim (\mathbf{M}_{2} + \mathbf{M}_{1}^{0}) = \dim \mathbf{E} + \dim (\mathbf{M}_{1} \cap \mathbf{E}^{0})\) (consider the canonical map from E to \(E/E^{0}\) ).
\end{enumerate}

\begin{enumerate}
\def\labelenumi{\arabic{enumi})}
\setcounter{enumi}{2}
\tightlist
\item
  Let K be a commutative ring, \(f\) a monic polynomial in K{[}X{]}, of degree \(n \geqslant 1\) ; let A be the quotient algebra K{[}X{]}/(f) admitting as a basis over K the elements 1, \(\xi, \xi^2, \ldots, \xi^{n-1}\) (chap.~IV, § 1, no. 5, prop. 4). Show that the data of an A-module E is equivalent to the data of a K-module \(E_0\) and of a K-endomorphism \(j\) of \(E_0\) such that \(f(j) = 0\) . For every
\end{enumerate}

\(u \in \mathrm{E}^*\) , set \(u(x) = \sum_{k=0} u_k(x)\xi^k\) ; show that if \(\alpha_0 = f(0)\) is invertible in K, the map \(u \to u_0\) is a K-isomorphism of \(\mathbf{E}^*\) on \((\mathbf{E}_0)^*\) whose reciprocal isomorphism will be made explicit.

¶ 4) a) Let A be a ring (commutative or not), σ an automorphism of A such that there exists an invertible element γ ∈ A satisfying γ\^{}σ = γ, and such that one has ξ\^{}σ² = γξγ\^{}-1 for all ξ ∈ A. Let B be a left A-module having a basis of two elements (e₁, e₂); show that one defines on B a ring structure by taking as multiplication in B the composition law

\[
(\xi e _ {1} + \eta e _ {2}) (\xi^ {\prime} e _ {1} + \eta^ {\prime} e _ {2}) = (\xi \xi^ {\prime} + \eta \eta^ {\prime \sigma} \gamma) e _ {1} + (\eta \xi^ {\prime \sigma} + \xi \eta^ {\prime}) e _ {2}.
\]

For this ring structure, \(e_1\) is the identity element (identified with the identity element 1 of A); if we set \(e_2 = \rho\) , one has \(\rho^2 = \gamma\) and \(\rho\xi = \xi^\sigma\rho\) for all \(\xi \in A\) . Moreover, B is a right A-module, of which 1 and \(\rho\) form a basis. If A is a field, a necessary and sufficient condition for B to be a field is that \(\gamma\) is not of the form \(\lambda^\sigma\lambda\) (where \(\lambda \in A\) ). (Cf. Chap. VIII, § 12, Exerc. 8).

\begin{enumerate}
\def\labelenumi{\alph{enumi})}
\setcounter{enumi}{1}
\tightlist
\item
  Let J be an involutive antiautomorphism of A. Suppose there exists an invertible element \(\delta\in\mathsf{A}\) satisfying the following conditions:
\end{enumerate}

\[
\delta^ {J} = \delta , \quad \delta^ {\sigma} \delta = \gamma \gamma^ {J}, \quad (\xi^ {J}) ^ {\sigma} = \delta (\xi^ {\sigma}) ^ {J} \delta^ {- 1} \quad \text { for   all } \quad \xi \in A.\tag{1}
\]

Show that one can then extend J to an involutive antiautomorphism of B (still denoted J) by setting

\[
(\xi + \eta \rho) ^ {J} = \xi^ {J} + \gamma^ {- 1} \delta^ {\sigma} (\eta^ {J}) ^ {\sigma} \rho .\tag{2}
\]

If in addition A and B are fields, and if \(\sigma\) is not an inner automorphism, show that conditions (1) are necessary for the existence of an extension of J to an involutive antiautomorphism of B (by setting \(\rho^{J} = \alpha + \beta\rho\) , with \(\alpha, \beta\) in A, show that one necessarily has \(\alpha = 0\) by writing the condition \((\rho\xi)^{J} = \xi^{J}\rho^{J}\) for all \(\xi \in A\) ).

\begin{enumerate}
\def\labelenumi{\alph{enumi})}
\setcounter{enumi}{2}
\tightlist
\item
  We assume that conditions (1) are satisfied and that the involutive antiautomorphism J is extended to B by (2). Let F be a unitary B-module, E the unitary A-module obtained from F by restricting the ring of scalars to A; the map \(j: x \to \rho x\) is then a bijective semilinear map from E onto itself, relative to the automorphism \(\sigma\) of A and such that \(j^{2}(x) = \gamma x\) . Let \(\Phi\) be a Hermitian form on F (for J); for \(x \in E\) , \(y \in E\) , set \(\Phi(x, y) = \Phi_{1}(x, y) + \Phi_{2}(x, y)\rho\) , where \(\Phi_{1}(x, y) \in A\) , \(\Phi_{2}(x, y) \in A\) . Show that \(\Phi_{1}\) is a Hermitian form on E (with respect to J), such that
\end{enumerate}

\[
\Phi_ {1} (j (x), j (y)) = (\Phi_ {1} (x, y)) ^ {\sigma} \delta ;
\]

we have \(\Phi_{2}(x,y)=\Phi_{1}(x,j(y))\delta^{-1}\) , \(\Phi_{2}\) is a sesquilinear form on E for the antiautomorphism (in general non-involutive) \(\xi\to(\xi^{\mathrm{J}})^{\sigma}\) , such that

\[
\Phi_ {2} (j (x), j (y)) = (\Phi_ {2} (x, y)) ^ {\sigma} \delta^ {\sigma}
\]

and

\[
\Phi_ {2} (y, x) = \gamma^ {\mathrm{J}} \delta^ {- 1} ((\Phi_ {2} (x, y)) ^ {\mathrm{J}}) ^ {\sigma}.
\]

Converses. For \(\Phi\) to be nondegenerate, it is necessary and sufficient that \(\Phi_{1}\) or \(\Phi_{2}\) be nondegenerate. Special case where B is a quaternion algebra over a commutative ring K, corresponding to a pair \((\alpha, \beta)\) of elements of K, \(\beta\) being invertible, A the subalgebra K + Ku of B, and J and \(\sigma\) the map \(\xi \to \bar{\xi}\) (chap.~II, § 7, no. 8).

\begin{enumerate}
\def\labelenumi{\alph{enumi})}
\setcounter{enumi}{3}
\tightlist
\item
  Let u be an automorphism of the A-module E. For u to be an automorphism of the B-module F, leaving the form \(\Phi\) invariant, it is necessary and sufficient that u satisfy any two of the following three conditions:
\end{enumerate}

\(1^{\circ} u\) leaves \(\Phi_{1}\) invariant;

\(2^{o} u\) leaves \(\Phi_{2}\) invariant;

\(3^{o} u\) commutes with \(j\) .

\begin{enumerate}
\def\labelenumi{\arabic{enumi})}
\setcounter{enumi}{4}
\tightlist
\item
  Let A be a commutative ring, E an A-module, \((x_{i})_{i\in I}\) a system of generators of E; let R be the submodule of the module \(\mathrm{A}^{(\mathrm{I})}\) formed by the elements \((y_{i})_{i\in\mathrm{I}}\) such that \(\sum_{i}y_{i}x_{i}=0\) , and let \((a_{\lambda})_{\lambda\in\mathrm{L}}\) be a system of generators of R (with \(a_{\lambda}=(a_{\lambda i})_{i\in\mathrm{I}}\) ). Let \((b_{ij})\) be a family of elements of A \((i\in\mathrm{I},j\in\mathrm{I})\) . For there to exist a quadratic form Q such that \(\mathrm{Q}(x_{i})=b_{ii}\) and \(\Phi(x_{i},x_{j})=b_{ij}\) for \(i\neq j\) ( \(\Phi\) denoting the bilinear form associated with Q), it is necessary and sufficient that \(b_{ij}=b_{ji}\) for all i, j in I, and that, for all \(\lambda\in L\) and \(i\in I\) one has
\end{enumerate}

\[
\sum_ {\{i, j \}} b _ {i j} a _ {\lambda i} a _ {\lambda j} = \sum_ {j \neq i} b _ {i j} a _ {\lambda j} + 2 b _ {i i} a _ {\lambda i} = 0;\tag{3}
\]

we then have \(\mathrm{Q}(\sum_{i}a_{i}x_{i})=\sum_{\{i,j\}}b_{ij}a_{i}a_{j}\) . Deduce from this a new proof of prop. 3 of no. 4. (Note that the \(x_{i}^{\prime}=1\otimes x_{i}\) form a system of generators of \(A^{\prime}\otimes_{A}E\) , and that the \(A^{\prime}\) -module \(A^{\prime}\otimes_{A}E\) is isomorphic to \(A^{\prime(I)}/R^{\prime}\) , where \(A^{\prime(I)}\) is identified with \(A^{\prime}\otimes_{A}A^{(I)}\) and \(R^{\prime}\) is generated by the image of R under the canonical map from \(A^{(I)}\) into \(A^{\prime(I)}\) ).

\begin{enumerate}
\def\labelenumi{\arabic{enumi})}
\setcounter{enumi}{5}
\tightlist
\item
  Let A be a commutative ring of characteristic 2, E a free A-module, \(\mathcal{A}\) (resp. \(\mathcal{S},\mathcal{Q}\) ) the A-module of alternating (resp. symmetric bilinear, quadratic) forms on E. We have \(\mathcal{A}\subset\mathcal{S}\) ; one also defines a linear map \(\omega\) of \(\mathcal{S}\) in \(\mathcal{Q}\) , and a linear map \(\theta\) of \(\mathcal{Q}\) in \(\mathcal{A}\) as follows: for every bilinear form \(\Phi\in\mathcal{S}\) , \(\omega(\Phi)\) is the quadratic form \(x\to\Phi(x,x)\) , and for every quadratic form \(Q\in\mathcal{Q}\) , \(\theta(Q)\) is the bilinear form associated with Q, which is alternating. Show that \(\omega(0)=\mathcal{A}\) , \(\theta(\mathcal{Q})=\mathcal{A}\) and \(\bar{\theta}(0)=\omega(\mathcal{S})\) .
\end{enumerate}

¶ 7) Let A be a commutative ring, and E, F two A-modules. A map Q from E to F is said to be quadratic if it satisfies the following conditions: \(1^{\circ} \mathrm{Q}(\alpha x) = \alpha^{2} \mathrm{Q}(x)\) for \(\alpha \in \mathrm{A}\) , \(x \in \mathrm{E}\) ; \(2^{\circ}\) the map \((x, y) \to \mathrm{Q}(x + y) - \mathrm{Q}(x) - \mathrm{Q}(y)\) of \(\mathrm{E} \times \mathrm{E}\) into F is bilinear. If \(f\) is a linear map from an A-module \(\mathrm{E}_{1}\) into E, \(\mathrm{Q} \circ f\) is a quadratic map from \(\mathrm{E}_{1}\) into F.

\begin{enumerate}
\def\labelenumi{\alph{enumi})}
\tightlist
\item
  Let E be an A-module, \(\mathrm{A}^{(\mathrm{E})}\) the module of formal linear combinations of elements of E with coefficients in A (chap. II, § 1, no. 8), and for every \(x \in \mathrm{E}\) , let \(\varepsilon_x\) be the corresponding element of the standard basis of \(\mathrm{A}^{(\mathrm{E})}\) . Let \(\Gamma^2(\mathrm{E})\) be the quotient of \(\mathrm{A}^{(\mathrm{E})} \times (\mathrm{E} \otimes_{\mathbf{A}} \mathrm{E})\) by the submodule R generated by the elements ( \(\varepsilon_{x+y} - \varepsilon_x - \varepsilon_y, -x \otimes y\) ) and ( \(\varepsilon_{\lambda x} - \lambda^2\varepsilon_x, 0\) ), for \(x \in \mathrm{E}\) , \(y \in \mathrm{E}\) , \(\lambda \in \mathrm{A}\) . For every \(x \in \mathrm{E}\) , set \(\gamma(x) = \varphi(\varepsilon_x, 0)\) , denoting by \(\varphi\) the canonical mapping from \(\mathrm{A}^{(\mathrm{E})} \times (\mathrm{E} \otimes \mathrm{E})\) on \(\Gamma^2(\mathrm{E})\) ; we say that \(\gamma\) is the canonical mapping from E into \(\Gamma^2(\mathrm{E})\) . Show that \(\gamma\) is a quadratic map from E to \(\Gamma^2(\mathrm{E})\) and that, for every quadratic map Q from E into an A-module F, there exists one and only one linear map \(q\) of \(\Gamma^2(\mathrm{E})\) in F such that \(Q = q \circ \gamma\) (in other words, \((\Gamma^2(\mathrm{E}), \gamma)\) is a solution of a universal mapping problem; cf.~Sets, chap.~IV, § 3).
\end{enumerate}

For every pair of A-modules E, E′ and every linear map f from E into E′, show that, if γ′ denotes the canonical map from E′ into Γ²(E′), there exists one and only one linear map \(\bar{f}\) from Γ²(E) to Γ²(E') such that γ'∘f = \(\bar{f} \circ \gamma\) .

\begin{enumerate}
\def\labelenumi{\alph{enumi})}
\setcounter{enumi}{1}
\item
  Assume that E is the direct sum of two submodules M, N. Define a canonical isomorphism of \(\Gamma^{2}(\mathrm{E})\) onto the direct sum of modules \(\Gamma^{2}(\mathrm{M})\) , \(\Gamma^{2}(\mathrm{N})\) and \(\mathbf{M} \otimes \mathbf{N}\) (show that this direct sum is a solution of the same universal mapping problem as \(\Gamma^{2}(\mathrm{E})\) ).
\item
  Let F be a submodule of E, j the canonical injection of F into E. Define a canonical isomorphism of \(\Gamma^{2}(\mathrm{E}/\mathrm{F})\) onto
\end{enumerate}

\[
\Gamma^ {2} (\mathrm{E}) / (\bar {j} (\Gamma^ {2} (\mathrm{F})) + \psi (\mathrm{E} \times \mathrm{F})),
\]

where \(\psi(x, y) = \varphi(0, x \otimes j(y))\) for \(x \in \mathrm{E}, y \in \mathrm{F}\) . (Same method).

\begin{enumerate}
\def\labelenumi{\alph{enumi})}
\setcounter{enumi}{3}
\item
  Let A' be a commutative ring, and h a homomorphism from A into A'. Define a canonical isomorphism of \(\Gamma^{2}(\mathbf{A}'\otimes_{\mathbf{A}}\mathbf{E})\) onto \(\mathbf{A}'\otimes_{\mathbf{A}}\symbf{\Gamma}^{2}(\mathbf{E})\) (same method).
\item
  There exists one and only one linear map \(s\) of \(\Gamma^{2}(\mathrm{E})\) into \(\mathrm{E}\otimes\mathrm{E}\) such that \(s(\gamma(x))=x\otimes x\) for all \(x\in\mathrm{E}\) ; show that if \(\mathrm{E}\) is a free module, s is an isomorphism onto the submodule of symmetric tensors of order 2 on E.
\item
  Assume that A = Z and that E is a finite cyclic group of order n. Show that \(\Gamma^{2}(E)\) is a cyclic group of order n if n is odd, and of order 2n if n is even. (First note that if a is a generator of E, \(\gamma(a)\) is a generator of \(\Gamma^{2}(E)\) , and that \(\gamma(-ha) = \gamma(ha)\) for every integer h; deduce from this that if n is odd, \(n\gamma(a) = 0\) by taking \(h = (n - 1)/2\) ; similarly show that \(2n\gamma(a) = 0\) if n is even. Finally, prove that if n is odd (resp. even), there exists a quadratic map Q from E into a cyclic group of order n (resp. 2n) sending a onto a generator of this group.
\end{enumerate}

\begin{enumerate}
\def\labelenumi{\arabic{enumi})}
\setcounter{enumi}{7}
\tightlist
\item
  Let A be a commutative field, and let E, F be two vector spaces over A. Let g be a map from E to F, such that there exist three maps a, b, c from E × E to F, satisfying the identity
\end{enumerate}

\[
g (\lambda x + \mu y) = \lambda^ {2} a (x, y) + \lambda \mu b (x, y) + \mu^ {2} c (x, y)\tag{4}
\]

for all x, y in E, λ, μ in A.

\begin{enumerate}
\def\labelenumi{\alph{enumi})}
\tightlist
\item
  Show that we have \(a(x, y) = g(x)\) , \(c(x, y) = g(y)\) , \(b(y, x) = b(x, y)\) and \(b(\lambda x, y) = \lambda b(x, y)\) ; moreover, if we set
\end{enumerate}

\[
d (x, y, z) = b (x + y, z) - b (x, z) - b (y, z)
\]

show that we have \(d(x,y,z)=d(y,z,x)=d(z,x,y)\) and conclude from this that \(d(\lambda x,\mu y,\nu z)=\lambda\mu\nu d(x,y,z)\) .

\begin{enumerate}
\def\labelenumi{\alph{enumi})}
\setcounter{enumi}{1}
\tightlist
\item
  From a), deduce that if \(\mathbf{A} \neq \mathbf{F}_2\) one necessarily has \(d(x, y, z) = 0\) and consequently that \(g\) is a quadratic map (exerc. 7). On the contrary, if \(\mathbf{A} = \mathbf{F}_2\) and if dim \(\mathrm{E} \geqslant 3\) show that there exist maps \(g\) from E to A, satisfying (4) and for which \(d(x, y, z)\) is not identically zero.
\end{enumerate}

\begin{enumerate}
\def\labelenumi{\arabic{enumi})}
\setcounter{enumi}{8}
\tightlist
\item
  Let A be a commutative ring, E an A-module having a basis of n elements, Φ a symmetric bilinear form on E. Let e be an element of \(\bigwedge^{n}\) E forming a basis of this module, \(\Delta\) the discriminant of \(\Phi\) with respect to \(e\) , \(\varphi_{p}\) the canonical isomorphism of \(\bigwedge^{p}\) E onto \(\bigwedge^{n-p}\) E* relative to \(e\) (chap.~III, §8, no. 5). For \(x \in \bigwedge^{p}\) E, let \(d_{(p)}(x) = \bar{\varphi}_{n-p}^{-1}(d_{\Phi(p)}(x))\) ; show that the map \(d_{(p)}\) of \(\bigwedge^{p}\) E into \(\bigwedge^{n-p}\) E possesses the following properties:
\end{enumerate}

\begin{enumerate}
\def\labelenumi{\alph{enumi})}
\item
  For every \(x \in \wedge \mathrm{E}\) , \(d_{(n - p)}(d_{(p)}(x)) = (-1)^{p(n - p)}\Delta x\) .
\item
  For every pair of elements \(x, y\) of \(\wedge\) E, we have
\end{enumerate}

\[
x \wedge d _ {(p)} (y) = \Phi_ {(p)} (x, y) e \quad \text { and } \quad \Phi_ {(n - p)} (d _ {(p)} (x), d _ {(p)} (y)) = \Delta \Phi_ {(p)} (x, y).
\]

\begin{enumerate}
\def\labelenumi{\alph{enumi})}
\setcounter{enumi}{2}
\item
  We further assume that A is a field and that \(\Phi\) is nondegenerate. Then, if \(x\) is a decomposable \(p\) -vector \(\neq 0\) corresponding to a subspace F of E (chap. III, § 7, no. 3), \(d_{(p)}(x)\) is a decomposable \((n-p)\) -vector \(\neq 0\) corresponding to the subspace F \(^{0}\) orthogonal to F.
\item
  Extend the previous results to the case where (A being a commutative ring), \(\Phi\) is a sesquilinear form \(\varepsilon\) -hermitian for an involutive automorphism \(J \neq 1\) of A.
\end{enumerate}

\begin{enumerate}
\def\labelenumi{\arabic{enumi})}
\setcounter{enumi}{9}
\tightlist
\item
  \begin{enumerate}
  \def\labelenumii{\alph{enumii})}
  \tightlist
  \item
    Let A be a commutative ring, E an A-module having a basis of 3 elements, \(\Phi\) a symmetric bilinear form on E. With the notations of exerc. 9, for any two elements \(x\) , \(y\) of E, one sets \(x \overline{\wedge} y = d_{(2)}(x \wedge y)\) , and one says that this element is the vector product
  \end{enumerate}
\end{enumerate}

\[
\stackrel {{3}} {\wedge} \mathrm{E})
\]

of \(x\) and of \(y\) (relative to \(\Phi\) and to the basis \(e\) of \(\wedge E\) ). Show that \((x, y) \to x \overline{\wedge} y\) is an alternating bilinear map from \(E \times E\) into \(E\) , and that \(x \overline{\wedge} y\) is orthogonal to \(x\) and to \(y\) .

\begin{enumerate}
\def\labelenumi{\alph{enumi})}
\setcounter{enumi}{1}
\tightlist
\item
  Let \(\alpha, \beta\) be two invertible elements of A, B the quaternion algebra over A corresponding to the pair \((\alpha, \beta)\) (chap.~II, § 7, no. 8), 1, u, v, w the canonical basis of B over A; let E be the submodule of B having u, v, w as basis. Show that if x, y are two quaternions belonging to E, one has
\end{enumerate}

\[
x y = \Phi (x, y) + x \overline {{{{\wedge}}}} y
\]

where \(\Phi\) is a symmetric bilinear form on E, such that the linear maps associated with \(\Phi\) are bijective, and \(x\overline{\wedge}y\) is the vector product of x and y relative to the form \(\Phi\) and to the basis \(\alpha^{-1}\beta^{-1}u\wedge v\wedge w\) of \(\bigwedge^{3}\mathbf{E}\) .

\begin{enumerate}
\def\labelenumi{\arabic{enumi})}
\setcounter{enumi}{10}
\tightlist
\item
  Let \(\Phi\) be a nondegenerate \(\varepsilon\) -hermitian sesquilinear form on a finite-dimensional vector space E. One says that a vector subspace M of E is weakly orthogonal to a vector subspace N (relative to \(\Phi\) ) if one of the two subspaces M, N° contains the other.
\end{enumerate}

\begin{enumerate}
\def\labelenumi{\alph{enumi})}
\item
  Show that the relation “M is weakly orthogonal to N” is symmetric.
\item
  If M and N are weakly orthogonal, \(M^{0}\) and \(N^{0}\) are weakly orthogonal.
\item
  If M and N are weakly orthogonal and if M ∩ N = \{0\}, M and N are orthogonal.
\end{enumerate}

\section{Totally isotropic subspaces. Witt's theorem}

In this section one assumes, unless expressly stated otherwise, that A is a field. One denotes by \(\Phi\) , either a form \(\varepsilon\) -hermitian on E (with respect to the involutive antiautomorphism \(\lambda \to \overline{\lambda}\) of A), or the symmetric bilinear form associated with a quadratic form Q on E (A being assumed commutative in the latter case).

\subsection{Isotropic subspaces.}

DEFINITION 1. --- Given a module E over the ring A, an element x of E is said to be isotropic if \(\Phi(x, x) = 0\) . A submodule F of E is said to be

\begin{enumerate}
\def\labelenumi{\arabic{enumi})}
\item
  isotropic if there exists an element \(x \neq 0\) of F orthogonal to F;
\item
  totally isotropic if the restriction of \(\Phi\) to F is zero.
\end{enumerate}

When the module E is equipped with a quadratic form Q, an element of E will be said to be isotropic (resp. a submodule of E will be said to be isotropic, or totally isotropic) if this element is isotropic (resp. if this submodule is isotropic, or totally isotropic) with respect to the bilinear form associated with Q.

An isotropic vector is none other than a vector orthogonal to itself. To say that a submodule F is isotropic means that \(F \cap F^{0} \neq \{0\}\) , or equivalently that the restriction of \(\Phi\) to F is degenerate; a non-isotropic submodule G of E is therefore a submodule such that the restriction of \(\Phi\) to G is nondegenerate. For a submodule F of E to be totally isotropic, it is necessary and sufficient that one has \(F \subset F^{0}\) . If F is a totally isotropic submodule of E, the same is true of every submodule \(F'\) contained in F. The sum of a family of totally isotropic submodules that are pairwise orthogonal is a totally isotropic submodule. The set of totally isotropic submodules of E, ordered by inclusion, is evidently inductive; it follows that every totally isotropic submodule is contained in a maximal totally isotropic submodule.

PROPOSITION 1. --- Suppose that A is a field. Let F be a finite-dimensional non-isotropic subspace of E; then E is the direct sum of F and F⁰.

Indeed, since the restriction \(\Phi'\) of \(\Phi\) to F is nondegenerate by hypothesis, the map \(d_{\Phi'}\) from F to its dual F* associated on the right to \(\Phi'\) is injective, hence bijective since F and F* are two spaces of the same finite dimension. Consequently, for every \(y \in E\) , there exists one and only one element \(y_0\) of F such that one has \(\Phi(x, y) = \Phi(x, y_0)\) for all \(x \in F\) , that is to say \(y - y_0 \in F^0\) ; this proves that E is the direct sum of F and F \(^0\) .

COROLLARY. --- If F is a finite-dimensional subspace of E, and if Φ is nondegenerate, the following conditions are equivalent:

\begin{enumerate}
\def\labelenumi{\alph{enumi})}
\item
  F is non-isotropic.
\item
  \(\mathbf{F}^0\) is non-isotropic.
\item
  E is the direct sum of F and F \(^{0}\) .
\end{enumerate}

Proposition 1 indeed shows that \(a)\) implies \(c)\) , and \(c)\) implies \(a)\) and \(b)\) . Finally, if \(\mathbf{F}^0\) is non-isotropic, we have \(\mathbf{F} \cap \mathbf{F}^0 \subset \mathbf{F}^0 \cap \mathbf{F}^{00} = \{0\}\) ; hence \(\mathbf{F}\) is non-isotropic, which shows that \(b)\) implies \(a)\) .

DEFINITION 2. --- Let Q be a quadratic form on E. An element x of E is said to be singular (relative to Q) if Q(x) = 0. A submodule F of E is said to be:

\begin{enumerate}
\def\labelenumi{\arabic{enumi})}
\item
  singular if there exists an element \(x \neq 0\) of F which is singular and orthogonal to F;
\item
  totally singular if the restriction of Q to F is zero.
\end{enumerate}

The kernel of the quadratic module (E, Q) (§ 3, no. 4) consists of the singular elements of E°; for a submodule F to be singular, it is necessary and sufficient that its kernel be ≠ \{0\} . As Φ(x, y) = Q(x + y) − Q(x) − Q(y), every totally singular submodule ≠ \{0\} is singular. Since Φ(x, x) = 2Q(x), every singular vector is isotropic and every singular (resp. totally singular) submodule is isotropic (resp. totally isotropic); the converse holds if A is a field of characteristic ≠ 2. Every submodule contained in a totally singular submodule is itself totally singular. The sum of a family of totally singular submodules that are pairwise orthogonal is a totally singular submodule. The set of totally singular submodules of E, ordered by inclusion, is inductive; hence every totally singular submodule of E is contained in a maximal totally singular submodule.

\subsection{Witt decomposition.}

In addition to the conventions already in force since the beginning of the present paragraph, we shall add the following:

Condition (T). --- For all \(x \in E\) , there exists \(\alpha \in A\) such that \(\Phi(x, x) = \alpha + \varepsilon\bar{\alpha}\) .

This condition is always satisfied when \(\Phi\) is alternating, or when \(\varepsilon = 1\) and A is a field of characteristic \(\neq 2\) , taking then \(\alpha = \frac{1}{2}\Phi(x, x)\) (cf. Exercises 1 and 14).

Lemma 1. --- Let \(\Phi\) be a form \(\varepsilon\) -hermitian satisfying (T) (resp. the bilinear form associated with a quadratic form Q) on E, and let F be a totally isotropic (resp. totally singular) subspace of E, not reduced to 0. For every \(x \in E\) not orthogonal to F and all \(\alpha \in A\) , there exists \(y \in F\) such that

\[
\Phi (x + y, x + y) = \alpha + \varepsilon \bar {\alpha} \quad (\text { resp. } Q (x + y) = \alpha).
\]

Indeed, let us set \(\Phi(x, x) = \beta + \varepsilon\bar{\beta}\) (resp. \(Q(x) = \beta\) ). For \(y \in F\) we then have \(\Phi(x + y, x + y) = (\beta + \Phi(x, y)) + \varepsilon(\overline{\beta + \Phi(x, y)})\) since \(\Phi(y, y) = 0\) (resp. \(Q(x + y) = \beta + \Phi(x, y)\) since \(Q(y) = 0\) ). Since \(x\) is not orthogonal to \(F\) , the affine linear function \(y \to \Phi(x, y) + \beta\) on \(F\) is not constant; it therefore takes the value \(\alpha\) for some element \(y\) of \(F\) , which thus answers the question.

A Witt decomposition of E is any decomposition of E as a direct sum of three subspaces F, F', G such that F and F' are totally isotropic (resp. totally singular) and G is non-isotropic and orthogonal to F + F'; if E is finite-dimensional, the matrix of Φ with respect to a basis of E adapted to a Witt decomposition of E takes the form

\[
\left( \begin{array}{c c c} 0 & U & 0 \\ \varepsilon \overline {{U}} & 0 & 0 \\ 0 & 0 & V \end{array} \right)\tag{1}
\]

We say that \(\Phi\) is a neutral form if it is nondegenerate and if E is finite-dimensional and is the direct sum of two totally isotropic (resp. totally singular) subspaces. The direct sum of two neutral forms is a neutral form.

PROPOSITION 2. --- Let Φ be a nondegenerate ε-hermitian form satisfying (T) (resp. the bilinear form associated with a nondegenerate quadratic form Q), and let F be a totally isotropic (resp. totally singular) subspace of finite dimension r.

\begin{enumerate}
\def\labelenumi{\alph{enumi})}
\item
  If \(F'\) is a totally isotropic subspace of dimension r such that \(F' \cap F^{0} = \{0\}\) , then \(F + F'\) is non-isotropic and, for any basis \((f_{i})\) of F, there exists a basis \((f_{i}')\) of \(F'\) such that \(\Phi(f_{i}, f_{j}') = \delta_{ij}\) (Kronecker index) for \(i, j = 1, \ldots, r\) .
\item
  If G is a totally isotropic (resp. totally singular) subspace of dimension \(\leqslant r\) such that \(G \cap F^{0} = \{0\}\) , there exists a totally isotropic (resp. totally singular) subspace \(F' \supset G\) of dimension r such that \(F' \cap F^{0} = \{0\}\) .
\end{enumerate}

Let \(\Psi\) be the restriction of \(\Phi\) to \(F \times F'\) ; for \(x' \in F'\) , the relation \(\langle \Phi(x, x') = 0\) for all \(x \in F\rangle\) implies \(x = 0\) since \(F' \cap F^0 = \{0\}\) . The assertion \(a)\) then follows from the corollary of Proposition 6 of §1, no. 6, except for the fact that \(F + F'\) is non-isotropic. Now the subspace \(H = (F + F') \cap (F + F')^0\) is equal to \((F + F') \cap F^0 \cap F'^0\) . Since \(F \subset F^0\) , one has \((F + F') \cap F^0 = F + (F' \cap F^0) = F\) , whence \(H = F \cap F'^0\) ; hence \(H = \{0\}\) since we have seen that \(\Psi\) is nondegenerate. This indeed proves that \(F + F'\) is non-isotropic.

To demonstrate \(b\) ), we proceed by descending induction on \(s = \dim G\) . It thus suffices for us to prove that, if \(s < r\) , there exists a totally isotropic (resp. totally singular) subspace \(G'\) containing \(G\) , of dimension \(s + 1\) , and such that \(G' \cap F^0 = \{0\}\) . Since \(\dim G < \dim F\) , the restriction of \(\Phi\) to \(F \times G\) is degenerate, and since \(G \cap F^0\) is zero, \(F \cap G^0\) is nonzero. If one then had \(G + F^0 \supset G^0\) one would deduce, by taking the orthogonal subspaces and noting that \(F = F^{00}\) and \(G = G^{00}\) ( \(\S 1\) , no. 6, cor. 1 of prop. 4), that \(G^0 \cap F \subset G\) , whence

\[
\mathrm{G} ^ {0} \cap \mathrm{F} \subset \mathrm{G} \cap \mathrm{F} \subset \mathrm{G} \cap \mathrm{F} ^ {0} = \{0 \},
\]

which is impossible. There then exists an element \(x\) of \(G^0\) such that \(x \notin G + F^0\) ; since \(F \subset F^0\) , one can add to \(x\) a vector of \(G^0 \cap F\) without modifying these properties; since \(G^0 \cap F\) is totally isotropic (resp. totally singular) and \(\neq \{0\}\) , lemma 1 shows that one can choose \(x\) isotropic (resp. singular).

Then the subspace \(G' = G + Ax\) is of dimension \(s + 1\) , and is totally isotropic (resp. totally singular); moreover we have \(G' \cap F^{0} = \{0\}\) because, if \(y = z + ax\) ( \(z \in G, a \in A\) ) is in \(F^{0}\) , we have a = 0, since otherwise \(x \in F^{0} + G\) contrary to the choice of x, whence \(y = z \in G \cap F^{0} = \{0\}\) and y = 0. Consequently the subspace \(G'\) answers the question.

COROLLARY 1. --- If F is a totally isotropic (resp. totally singular) subspace of dimension r, there exists a totally isotropic (resp. totally singular) subspace F' of dimension r such that F ∩ F' = \{0\} and such that F + F' is non-isotropic.

It suffices to take G = \{0\} in prop. 2, b).

COROLLARY 2. --- Two neutral ε-hermitian forms on spaces of the same dimension over A are equivalent.

Remark. --- Under the conditions of corollary 1, E is the direct sum of F + F' and of the orthogonal of F + F'. One thus has a Witt decomposition of E. By prop. 2, a), there exist bases of F and F' such that, in the matrix (1) of Φ, the block U is the identity matrix 1.

PROPOSITION 3. --- Let \(\Phi\) be a nondegenerate \(\varepsilon\) -hermitian form satisfying (T) (resp. the bilinear form associated with a nondegenerate quadratic form Q). Let \(F_{1}\) and \(F_{2}\) be two maximal totally isotropic (resp. totally singular) subspaces of E, one of the two being finite-dimensional. Set \(F = F_{1} \cap F_{2}\) . Let \(S_{i}\) ( \(i = 1, 2\) ) be a supplement of F in \(F_{i}\) ; set \(S = S_{1} + S_{2}\) . There then exist two subspaces G and H of E such that

\begin{enumerate}
\def\labelenumi{\alph{enumi})}
\item
  The subspaces G + F, S, and H are non-isotropic and pairwise orthogonal;
\item
  E is the direct sum of F, S, G, and H;
\item
  There is no nonzero isotropic (resp. singular) vector in H;
\item
  G is totally isotropic (resp. totally singular).
\end{enumerate}

Additionally, \(F_{1}\) and \(F_{2}\) are both finite-dimensional and we have \(\dim \mathrm{F}_1 = \dim \mathrm{F}_2, \dim \mathrm{G} = \dim \mathrm{F}, \dim \mathrm{S}_1 = \dim \mathrm{S}_2, \operatorname{codim} \mathrm{H} = 2 \dim \mathrm{F}_1.\)

First, note that if N is a maximal totally isotropic (resp. totally singular) subspace, then every isotropic (resp. singular) vector x orthogonal to N belongs to N, for otherwise \(N + Ax\) would contradict the maximality of N. Thus, if, for i = 1, 2, \(x_{i}\) is an isotropic (resp. singular) vector of \(F_{i}^{0}\) , one has \(x_{i} \in F_{i}\) On the other hand, if y is an element of \(S_{1}\) orthogonal to \(S_{2}\) it is orthogonal to \(F_{1}\) since \(F_{1}\) is totally isotropic, hence to F, and consequently to \(F_{2} = S_{2} + F\) . Since y is isotropic (resp. singular) and is orthogonal to \(F_{2}\) , one has

\[
y \in S _ {1} \cap F _ {2} = S _ {1} \cap F _ {1} \cap F _ {2} = S _ {1} \cap F = \{0 \}.
\]

We therefore have \(S_1 \cap S_2^0 = \{0\}\) , and likewise \(S_2 \cap S_1^0 = \{0\}\) . Since one of the two subspaces \(F_1, F_2\) , for example \(F_1\) , is finite-dimensional, \(S_1\) is finite-dimensional, hence \(S_1^0\) has finite codimension (§ 1, no. 6, cor. 1 of prop. 4), and consequently \(S_2\) is finite-dimensional since \(S_2 \cap S_1^0 = \{0\}\) ; moreover this shows that one has dim \(S_2 \leqslant \text{codim } S_1^0 = \text{dim } S_1\) ; likewise dim \(S_1 \leqslant \text{dim } S_2\) , whence dim \(S_1 = \text{dim } S_2\) . Prop. 2 a) then shows that \(S = S_1 + S_2\) is non-isotropic.

This being so, the orthogonal N of S is non-isotropic (no. 1, cor. of prop. 1) and contains F; cor. 1 of prop. 2 therefore shows that there exists a subspace G totally isotropic (resp. totally singular) of N such that dim G = dim F, that G ∩ F = \{0\} and that G + F is non-isotropic. Thus d) is satisfied by G. One then satisfies a) and b) by taking for H the orthogonal of G + F in N. As for c), one notes that, since H is orthogonal to F₁ = S₁ + F, there is no nonzero isotropic (resp. singular) vector in H by virtue of what was remarked at the beginning of the proof and of the fact that H ∩ F₁ = \{0\}. Finally, certain of the assertions concerning dimensions were proved along the way; the others follow from them trivially.

COROLLARY 1. --- Under the hypotheses of prop. 3, two maximal totally isotropic (resp. totally singular) subspaces of finite dimension have the same dimension. For every maximal totally isotropic (resp. totally singular) subspace F of finite dimension, there exists another one \(\mathbf{F}'\) such that \(\mathbf{F} \cap \mathbf{F}' = \{0\}\) , and under these conditions \(\mathbf{F} + \mathbf{F}'\) is non-isotropic.

If \(F \cap F' = \{0\}\) , one will have \(G = \{0\}\) with the notations of prop. 3, and \(F + F'\) will be non-isotropic. The other assertions follow trivially from prop. 3 and cor. 1 of prop. 2.

COROLLARY 2. --- Let Q be a nondegenerate quadratic form on a vector space E of finite dimension n over an algebraically closed field A; there then exists a basis \((e_i)_{1 \leqslant i \leqslant n}\) of E such that

\begin{enumerate}
\def\labelenumi{(\arabic{enumi})}
\setcounter{enumi}{1}
\tightlist
\item
\end{enumerate}

\[
\mathrm{Q} (\sum_ {i = 1} ^ {n} x _ {i} e _ {i}) = \sum_ {i = 1} ^ {\nu} x _ {i} x _ {i + \nu} \quad \text {   if   } n = 2 \nu ,\tag{3}
\]

\[
\mathrm{Q} \left(\sum_ {i = 1} ^ {n} x _ {i} e _ {i}\right) = \sum_ {i = 1} ^ {\nu} x _ {i} x _ {i + \nu} + x _ {2 \nu + 1} ^ {2} \quad \text {   if   } n = 2 \nu + 1.
\]

Indeed, let \(F_{1}\) and \(F_{2}\) be two maximal totally singular subspaces such that \(F_{1} \cap F_{2} = \{0\}\) (cor. 1) and let q be their dimension. One then has \(G = \{0\}\) with the notations of prop. 3. Taking a basis \((e_{i})_{1 \leqslant i \leqslant q}\) of \(F_{1}\) and a basis \((e_{i})_{q+1 \leqslant i \leqslant 2q}\) of \(F_{2}\) such that \(\Phi(e_{i}, e_{j+q}) = \delta_{ij}\) for \(i, j = 1, \ldots, q\) (prop. 2 a)), we see that it suffices to show that dim \(H \leqslant 1\) . Now, if \(x \in H\) , \(y \in H\) and if \(x \neq 0\) , the equation \(Q(y - ax) = Q(y) - a\Phi(x, y) + a^{2}Q(x) = 0\) has at least one solution \(a_{0}\) since \(Q(x) \neq 0\) and one has \(y = a_{0}x\) since every singular vector of H is zero.

DEFINITION 3. --- Suppose that E is finite-dimensional and that \(\Phi\) is a nondegenerate \(\varepsilon\) -hermitian form satisfying (T) (resp. the bilinear form associated with a nondegenerate quadratic form Q). The index of \(\Phi\) (resp. Q) is the common dimension of the maximal totally isotropic (resp. totally singular) subspaces of E.

If n is the dimension of E and v the index of \(\Phi\) (resp. Q), Prop. 3 shows that one has

\[
n \geqslant 2 v.\tag{4}
\]

Moreover, since every totally isotropic (resp. totally singular) subspace is contained in a maximal totally isotropic (resp. totally singular) subspace, the maximal totally isotropic (resp. totally singular) subspaces are precisely those of dimension v. The assertion that \(\Phi\) (resp. Q) has index 0 means that every isotropic (resp. singular) vector of E is zero. In a space of even dimension \(n\) the neutral forms are those of index \(\frac{1}{2} n\) ; there is no neutral form in a space of odd dimension. Prop. 3 shows that every form is a direct sum of a neutral form and a form of index 0.

PROPOSITION 4. --- Let Q be a nondegenerate quadratic form on E such that there exists a vector \(x \neq 0\) of E with Q(x) = 0. For every element a of A, there then exists \(y \in E\) such that Q(y) = a.

Indeed, by cor. 1 of prop. 2, there exists a subspace \(\mathbf{G} = \mathbf{F} + \mathbf{F}'\) ( \(\mathbf{F}, \mathbf{F}'\) : totally singular subspaces of dimension 1) of \(\mathbf{E}\) of dimension 2, such that the restriction of \(\mathbf{Q}\) to \(\mathbf{G}\) is neutral. If \(\{e, e'\}\) ( \(e \in \mathbf{F}, e' \in \mathbf{F}'\) ) is a basis of \(\mathbf{G}\) , one has

\[
\mathrm{Q} (x e + x ^ {\prime} e ^ {\prime}) = b x x ^ {\prime} \quad (x \in \mathrm{A}, x ^ {\prime} \in \mathrm{A}, b \in \mathrm{A}, b \neq 0).
\]

It therefore suffices to take for y the vector \(ae + b^{-1}e'\) .
